% 本文件由 LaTeX 排版 Agent 根据有效 translation.json 自动生成。
% author=Athanasius; work_id=2801; title=驳异教徒

\chapter{\WorkTitle{驳异教徒}}

% structure: p0001 source=h1 target=omitted reason=page-title-already-rendered-by-parent

\Emph{引言：——本书的目的在于为基督宗教的教义，尤其是十字架辩护，反驳外邦人的嘲笑和反对。该教义的效果乃是其主要的证明。}\par

% structure: p0003 source=h2 target=section reason=relative-heading-rank; base=h2

\section{第一部分}

\Emph{引言：——本书的目的在于为基督宗教的教义，尤其是十字架辩护，反驳外邦人的嘲笑和反对。该教义的效果乃是其主要的证明。} 1. 我们宗教的知识与事物的真理，其本身是自明的，几乎无需人的教师，因为几乎日复一日，它藉着事实证实自己，并藉着基督的教义显明得比太阳更加光明。 2. 然而，\Person{玛卡里乌斯}，你仍渴望听到这些道理；那么，请来，容我们按我们所能，将基督信仰的几点内容陈明。你本有能力从圣经中寻得，却仍慷慨地愿意也听别人讲述。 3. 因为，虽然神圣而默感的\WorkTitle{圣经}已足以宣布真理，——同时我们真福的导师们为此而编撰的其它著作，若有人研读，就可对圣经的诠释获得一些知识，并能学到他所愿意知道的，——然而，既然我们手头目前没有导师们的著作，我们必须以书面形式将我们从他们那里所学到的，即救主基督的信仰，传达给你；以免有人轻视在我们当中所教导的道理，或认为对基督的信仰是不合理的。因为外邦人正是对此加以诽谤嘲笑，并大声讥讽我们，执拗于基督的十字架这一事实；然而，正是在这一点上，我们必须怜悯他们缺乏理智，因为他们当诽谤基督的十字架时，竟看不见它的能力充满了普世，且藉着它，对天主的认识的效果已向万民显明。 4. 因为，如果他们也真正留心到祂的天主性，就不会嘲笑这件事了。相反，他们自己也会认出这人就是世界的救主，十字架并非灾难，而是受造界的治愈。 5. 因为，如果十字架之后，一切偶像崇拜均被推翻，所有的魔鬼的显现都被这标记驱散，唯独基督受敬拜，且藉着祂认识圣父；同时，那些反对者蒙羞，祂每日无形地赢得了这些反对者的灵魂——那么，我们可以合理地反问他们：这件事如何还能被我们视为纯属人为，而不承认那位登上十字架的就是\OriginalTerm{圣言}{Word of God}、世界的救主呢？但在我看来，这些人就像一个人，在太阳被云遮蔽时诽谤太阳，却还惊叹于太阳的光芒，明明看到整个受造界都被它照亮。 6. 因为，就如光是高贵的，而作为光之首要原因的太阳则更高贵，同样，既然普世都充满了对祂的认识，这本身就是一件神圣的事，那么，如此成就的安排者与首要原因，就是天主和圣言。 7. 因此，我们按我们力所能及地陈述，首先驳斥不信者的无知；这样，谬误既被驳倒，真理便能自身彰显；而你自己，朋友，也能确信你所相信的是真实的，在认识基督上并没有受骗。此外，我认为对你这爱慕基督的人谈论基督是合宜的，因为我确信你对祂的信仰与认识，是高于其它一切事物的。\par

\Emph{恶并非事物本质本性的一部分。人在恩宠及对天主的认识中的最初创造与构造。}\par

起初，邪恶并不存在。即使在现今的圣人中，它也不存在，也绝不属于他们本性的任何部分。但后来人自己开始策划并制造它，自招损害。由此他们也发明了偶像，把不存在的东西当作存在。 2. 因为创造万有并为万王之王的这一天主，祂的实体超越一切本质和人的探索，祂既是善而极为高尚，便藉着祂自己的\OriginalTerm{圣言}{Word}、我们的救主耶稣基督，按照自己的肖像造成了人类，并使人在与祂自身相似的基础上，能够看见并认识现实；同时赐给人对祂永恒性的观念和认知，好使人保持其本性的完整，从不偏离对天主的观念，也不退避与圣者的共融；人因拥有那位施予者的恩宠，并从圣父的圣言获得天主自身的能力，就能喜乐并与天主性相交，度着一种不受伤害、真实蒙福的不死生命。因为人的理智对于认知天主性了无障碍，他藉其纯净，常常得见圣父的肖像——圣言天主，人自己就是按照祂的肖像而受造的。当他仰观通过\OriginalTerm{圣言}{Word}而普及宇宙的那个\OriginalTerm{天主眷顾}{Providence}时，他惊叹不已，超脱了感官的事物及一切形体的外在表现，却藉他心灵的能力，紧握天上的、属神的、可悟的事物。 3. 因为当人的理智不交结于形体，也不从外界混杂任何形体的贪欲，而是完全超乎它们之上，专务于自身，一如起初被造的状态时，那么，它便超越感官之事及一切人间事，被提升到高天之上；看见圣言，并在祂内也看见圣言的圣父，喜悦于瞻仰祂，并因朝向祂的渴望而获得更新； 4. 正如第一个受造的人——在希伯来文中被称为\Person{亚当}的——在\WorkTitle{圣经}中被描述为，起初他的理智在一种不因羞耻而局促的自由中朝向天主，并在那地方——圣\Person{梅瑟}用比喻称之为乐园的地方——与圣者们共享着对所悟之事的观望。因此，灵魂的纯净本身就足以反映天主，如主所说：\BibleQuote{心里洁净的人是有福的，因为他们要看见天主。}\BibleRef{玛 5:8}\par

\Emph{人因沉湎于物质事物而从上述境界堕落。}\par

因此，如我们所说，\OriginalTerm{造物主}{Creator}创造了人类，且原意就是要人类保持如此。但人类轻看更善美的事物，不愿领会它们，开始更热切地寻求那些更接近自己的事物。 2. 更接近人自身的，乃是\OriginalTerm{肉身}{body}及其\OriginalTerm{感官}{senses}；于是，他们将心思从由思想所领悟的现实中移开，转而注视自己；既如此行事，执着于肉身与其他感官之物，在自己周围的事物中彷佛受了欺骗，他们陷入了对自己的\OriginalTerm{贪欲}{lust}，宁取属己之物，而舍弃了对属于天主之事物的\OriginalTerm{默观}{contemplation}。随后，他们安居于这些事物之中，不愿离开如此临近他们的事物，将灵魂纠缠于肉身的\OriginalTerm{快乐}{pleasure}，被种种贪欲所折磨、搅扰，全然忘却了他们最初从天主所领受的\OriginalTerm{能力}{power}。 3. 这一真理，从最初受造的人身上，按圣经关于他的记载，便可得知。因为，他也一样，只要他保持心思专注于天主和对天主的默观，就会转离对肉身的关注。但当听从蛇的诱导，他离弃了对天主的顾念，开始注视自己时，他们不仅陷入肉身的贪欲，而且发现自己赤身露体，一发觉便感到羞耻。但他们知道自己赤身露体，并非指失去了衣服，乃是指他们已被剥夺了对神圣事物的默观，而将自己的理智转向了相反的事物。因为离弃了对那位唯一且真实者——即天主——的顾念，以及对祂的渴望，他们从此投入了各式各样的贪欲，以及各自身感官的贪欲之中。 4. 然后，正如事情发展所常见的那样，他们对每一件、每一物都产生了欲望，渐渐习惯了这些欲望，甚至不敢舍弃它们：由此，灵魂变得屈从于怯懦与惊恐，以及\OriginalTerm{有死}{mortality}的快乐与思想。因为不愿舍弃自己的贪欲，她害怕死亡，害怕与肉身分离。再者，由于贪欲未得满足，她学会了杀害和行不义。于是，我们自然要尽其所能地说明，她是如何如此行事的。\par

4. \Emph{灵魂因滥用\OriginalTerm{选择的自由}{freedom of choice}，而从真理逐渐坠落至谬误。}\par

灵魂既已离弃对理知事物的默观，尽情运用肉身的各种活动，且以对肉身的默观为乐，看出快乐对她而言是好的，便受了误导，妄用了善之名，以为快乐就是善的本质本身：就好比一个人精神失常，向所遇见的每一个人索要刀剑，却以为自己心智健全。 2. 然而，因爱上快乐，她开始以种种方式将其付诸实践。因为，她本性是\OriginalTerm{可动的}{mobile}，尽管已转离了善，却仍没有失去其\OriginalTerm{可动性}{mobility}。于是，她运动起来，但不再按照德行，也不再为了看见天主，而是幻想虚假的事物，对她的能力加以新奇的使用，将其滥用作获取她所虚构之快乐的手段，毕竟她受造时就具有支配自己的能力。 3. 因为她能够，一方面倾向于善，另一方面也能拒绝善；但在拒绝善的时候，她必然思及那与善相反者，因为她全然无法停止运动，正如我先前所说，她本性就是可动的。她既知自己拥有支配自身的能力，便看出自己能够将肉身的肢体用于两个方向：既可以面向那\Quote{存有}者，也可以面向那\Quote{非有}者。 4. 然而，善乃是\Quote{存有}，恶则是\Quote{非有}；因此，我所说的\Quote{存有}，指的是善，因为善以那自有者天主为典范。而我所说的\Quote{非有}，则指恶，因为恶在于人思想中的虚假幻象。因为，虽然肉身有眼睛，可以观看\OriginalTerm{受造界}{Creation}，并由其完全和谐的构造而认识造物主；有耳朵，听从神圣的谕令和天主的法律；有双手，既从事必要的劳作，又可以向天主举手祈祷；然而灵魂却离弃对善的默观，偏离其固有领域的运动，飘荡远离，转向相反之物。 5. 于是，正如我先前所说，她看见并滥用其能力，察觉到自己也将肉身的肢体反向移动：因此，她不注视受造界，反转眼目投向贪欲，显示出她也拥有这种能力；她以为只要在动，便是在维护自己的尊严，任凭己意而行也不犯罪；却不明白她被造不仅是为了运动，更是为朝向正确的方向运动。因此，宗徒有一句话向我们保证：\BibleQuote{凡事都可行，但不全有益}\BibleRef{格前 10:23}。\par

5. \Emph{因此，恶本质上在于拣选较低者而舍弃较高者。}\par

但人的狂妄，不顾及何谓相宜与合宜，只在意自身所能为，便反其道而行；由此，将手转向相反的方向，致使人犯下谋杀，又使听觉趋向悖逆，使其他肢体趋向淫乱，而非合法的生育；舌头不说正直的话，却去毁谤、辱骂和发虚誓；手则去偷窃和击打同胞；嗅觉趋向各种邪淫的气味；脚则迅速去流血，肚腹趋向醉酒和贪得无厌的饕餮。2. 这一切都是灵魂的恶习与罪：并非有任何别的原因，只是由于拒绝了更好的事物。就像一位驭车者，驾着赛车上了赛道，却不留意他应当驶向的目标，反而不顾目标，只是任意或随己意驱策马匹，时常撞向迎面而来的人，时常冲下陡坡，随着马车疾驰的速度随意狂奔，以为如此奔跑并未错过目标——因为他只顾奔跑，却未察觉自己已远远偏离目标——灵魂也是如此，一旦偏离那通向天主的道路，驱使身体的肢体超出所当行的，甚或由于自己的作为而随同肢体一同被驱使，便犯罪并自招灾祸，不晓得自己已经离开了正路，偏离了真理的目标；而那\OriginalTerm{怀有基督的人}{Christ-bearing man}，真福保禄，正是注视这目标而说：\BibleQuote{我向着目标直奔，为争取天主在基督耶稣内召我争夺的奖品。}\BibleRef{斐 3:14}因此，这圣洁的人，以善为目标，从未作恶事。\par

6. \Emph{关于恶之本性的错误观点：即认为恶是事物本性中的东西，并具有\OriginalTerm{实体存在}{substantive existence}。（a）异教思想家：（认为恶寓于物质之中）。对其的反驳。（b）异端教师：（\OriginalTerm{二元论}{Dualism}）。从圣经来的反驳。}\par

如今，某些希腊人偏离了正道，不认识基督，便赋予恶一种实体的独立存在。如此他们犯了双重错误：若恶具有自己独立的自立性与存在，他们要么否认造物主是万物的创造者；要么，若他们意指祂是万物之创造者，便必然要承认祂也是恶的制造者。因为依他们之见，恶被列入存在之物中。2. 但这必然显得荒谬而不可行。因为恶并非出于善，也不寓于善内，或源自善，否则善就不成其为善，因其本性已混有恶或成为恶的原因。3. 但那些背离教会教导，在\OriginalTerm{信德}{Faith}上遭受了船破之灾的\OriginalTerm{异端分子}{sectaries}（\BibleRef{弟前 1:19}），他们也错误地认为恶具有\OriginalTerm{实体存在}{substantive existence}。他们任意设想在\OriginalTerm{真天主}{the true One}，即我们主耶稣基督之父以外，另有一位神，且认为那神是恶的\OriginalTerm{非受造}{unmade}制造者、邪恶之首，也是创造宇宙的\OriginalTerm{工匠}{artificer}。然而，这些人不仅可藉由\WorkTitle{圣经}，亦可从他们这些疯狂幻想所由发的人的理智本身，轻易予以驳斥。4. 首先，我们的主及救主耶稣基督在祂自己的\WorkTitle{福音}中，确证梅瑟的话说：\BibleQuote{上主天主是唯一的；}又说：\BibleQuote{父啊，天地的主宰，我称谢你。}\BibleRef{谷 12:29}但若天主是唯一的，且同时是天地的主宰，那么在祂之外又怎能另有神呢？若唯一真天主充满天地之间的一切，那么他们所设想的神又有什么位置呢？按其救主所言，基督的天主父本身就是天地的主宰，又如何另有创造者呢？5. 除非他们真的说，这好比处于一种平衡状态，恶神能胜过善神。但若他们这样说，看他们堕落到何等不虔的地步。因为当力量相当时，优越者和较善者便无法辨明。因为若一方即使另一方不情愿也依然存在，那么双方就同等的强，也同等的弱：强，因为任何一方的存在都是对方意愿的失败；弱，因为所发生的事情与他们的意愿相反：因为尽管有恶神存在，善天主依然存在；同样，尽管有善天主存在，恶神也依然存在。\par

7. \Emph{从理性反驳\OriginalTerm{二元论}{dualism}：不可能有两位神。至于恶的真理，乃是教会所教导的：即恶源于并寓于幽暗灵魂的乖谬选择之中。}\par

更特别的是，他们遭遇以下反驳。如果可见之物是恶神的工作，那么善的天主的工作是什么呢？因为除了造物主\OriginalTerm{造物主}{Artificer}的工作之外，别无可睹。如果没有造物主的工作可以使人认识祂，那么又有什么证据表明善的天主确实存在呢？因为造物主乃是透过其作品而被认识的。 2. 或者，两个彼此相反的本原何以能存在？是什么使它们分离，以使之分开存在？因为它们不能同时存在，因为互相毁灭。但也不能互相包含，因为其本性不相混合，亦不相似。因此，那分离它们的显然将具有第三种本性，且其自身就是神。但这第三个事物可能是什么本性？善还是恶？这无法确定，因为它不能同时具有二者的本性。3. 可见，他们的这种想法显然已腐朽败坏，教会神学的真理则必显明：即恶从起初就不与天主同在，也不在天主之内，亦无任何实体的存在；而是人由于缺乏对美善的瞻视，开始随自己的喜好，为自己构想和幻想虚无之物。4. 正如有人当太阳照耀，整个大地被其光辉普照时，却紧闭双眼，在无黑暗处幻想黑暗，然后在仿佛的黑暗中行走徘徊，常常跌倒，滑下陡坡，以为黑暗而非光明——因为他幻想自己能看见，其实完全看不见——照样，人的灵魂也紧闭双眼，拒绝凭借它得以看见天主的能力，反而为自己幻想了恶，并在其中活动，却不知道她自以为行事，却一无所成。因为她所幻想的并非真实存在，她也没有持守自己原始的本性；她如今的状况显然是她自身混乱的产物。5. 因为她受造是为看见天主，并蒙祂光照；但她却随从己意，离弃天主，转而寻求可朽坏之物和黑暗，正如圣神在经中某处所说：\BibleQuote{天主造人原是正直，但人却发明了许多诡计。}\BibleRef{训 7:29}。因此，自起初，人就为自己发现、谋划和幻想了恶。但现在我们该说明他们如何陷入偶像崇拜的疯狂，好让你知道偶像的发明完全不在于善，而全在于恶。但凡起源于恶的，在任何方面都绝不能被称为善——它全然是恶。\par

八、\Emph{偶像崇拜的起源与此相似。灵魂因忘却天主而物化，沉湎于世俗事物，便将它们当作神明。人类堕入毫不希望的迷惑与迷信的深渊。}\par

如今，人类的灵魂不满足于发明恶，逐渐开始冒险行更坏的事。因为她经验到各种不同的快乐，被遗忘神圣事物的念头缠绕；更且贪恋和注视肉身的私欲，以及当下的事物和对此的看法，她便不再认为可见之物以外还有任何存在，或除短暂和肉体的事物外还有什么是善的；如此，她转身离开并遗忘自己原是依照善的天主的肖像所造，就不再凭藉内里的能力看见那她按祂肖像受造的天主圣言；反而离开自己，幻想到并不存在的东西。2. 因为她被肉身的私欲所纠缠，遮蔽了那如同在她内的镜子，她唯有借这镜子才能看见圣父的肖像，于是她不再看见灵魂所应注视之物，却被万事动摇，只看见感官接触的事物。因此，负载着所有肉欲，迷失在这些事物的印象中，她便幻想那被自己悟性遗忘了的天主，是存在于肉体和可感之物内的，遂将天主的名字赋予可见之物，只光荣那些她渴望且悦目的事物。3. 因此，恶便招致偶像崇拜；因为人既然学会了制造本非实有的恶，同样也为自己捏造了并无真实存在的神明。这就像有人潜入深海，不再看见光和由光显现之物，因为他的眼睛向下看，全身被水包围；他只感知深海中的事物，以为它们之外别无所有，认为他看见的才是唯一的真实；照样，古时的人丧失了理性，沉沦于肉欲和肉身的幻想中，遗忘了对天主的认识和光荣，他们的理智昏暗，或更可说随从了无理性，便用可见之物为自己制造神明，\BibleQuote{他们崇拜受造之物，而不是造物主}\BibleRef{罗 1:25}，将受造物奉为神，而不尊崇他们的主、天主、根源和造物主。4. 正如上述比喻，潜入深海的人越下潜就越进入更暗更深的地方，人类亦是如此。因为他们并不停留在简单的偶像崇拜形式中，也不持守起初的状态；他们在最初的状态中愈久，就发明愈多的新迷信；不满足于先前的恶，他们又充满其他的恶，进一步陷入极度的无耻，在不信中变本加厉。但神圣的圣经为此作证说：\BibleQuote{恶人陷入极深的邪恶时，便起藐视。}\BibleRef{箴 18:3}\par

九、\Emph{偶像崇拜的种种发展：对天体、元素、自然物、神话生物、人化的私欲、活人和死人的崇拜。\Person{安提诺乌斯}的事例，以及被神化的皇帝们。}\par

因为那时人类的理解力已离弃了天主；他们的思想和想象日趋低俗，便将应归于天主的尊荣，首先给予了苍天、太阳、月亮和星辰，认为它们不仅是神，而且还是比它们更低下的其他神祇的起因。随后，在他们更为阴暗的想象中，又进一步降格，将神的名号赋予了\OriginalTerm{上层的清气}{æther}、大气以及空中的事物。接下来，在邪恶的道路上更进一步，他们竟将构成物体的元素和原理，即热、冷、干、湿，奉为神明来崇拜。 2. 然而，正如那些完全跌倒的人像蜗牛一般在泥淖中爬行一样，人类中最不敬虔的那些人，从对天主的观念中越坠越低，随后竟将活人以及人的形象立为神，有些是活人，有些甚至在死后。此外，他们还谋划并想象出更恶的事，最终竟将天主那神圣而超越的名字，移用到石头、木块、陆地和水中的爬虫，以及无理性的野兽身上，将一切属神的尊荣归于它们，从而离弃了那位真实而唯一的天主，即基督的父。 3. 但愿这些愚顽之人的大胆妄为就止于此，不至于在亵渎的自我混乱中走得更远。因为有些人的理智已坠入如此深渊，心智陷入如此黑暗，竟为自己编造出根本不存在、在受造物中全无位置的神明。他们将理性与非理性混杂，将本性相异的事物揉合，便崇拜由此而生的产物为神，比如埃及人中的狗头、蛇头和驴头的神祇，以及利比亚人中的羊头\Person{阿蒙}。另一些人则把人体各部位分割开来，将头、肩、手、脚一一立为神并奉之为神，仿佛他们的宗教对完整无缺的整体身体仍不满足。 4. 还有些人，将不虔敬推至极致，竟把发明这些事物以及他们自身邪恶的动机——即快乐与情欲——奉若神明，并加以崇拜，例如他们的\Person{厄洛斯}，以及\Place{帕福斯}的\Person{阿佛洛狄忒}。另有一些人，仿佛在堕落方面互相竞争，竟敢将他们的统治者甚或其儿子立为神，或是出于对君主的尊崇，或是出于对其暴政的恐惧，例如他们崇敬的\Place{克里特}的\Person{宙斯}，以及\Place{阿卡迪亚}的\Person{赫耳墨斯}；还有\Place{印度}人中的\Person{狄奥尼索斯}，\Place{埃及}人中的\Person{伊西斯}、\Person{奥西里斯}和\Person{霍鲁斯}；以及我们时代的\Person{安提诺乌斯}，罗马皇帝\Person{哈德良}的宠臣。尽管人们知道他只是凡人，且非正人君子，反倒充满放荡，然而他们因惧怕那位下令崇拜他的人而崇拜他。因为\Person{哈德良}旅居\Place{埃及}时，他那供其享乐的\Person{安提诺乌斯}死后，便下令让人崇拜他；他本人在这青年死后仍恋慕着他，然而这一切都充分暴露了他自己，也为反对一切偶像崇拜提供了证据，表明偶像崇拜之在人间的出现，无非是出于幻想者的情欲。正如天主的智慧事先作证说：\BibleQuote{发明偶像乃是淫乱的开端。} \BibleRef{智 14:12} 5. 请不要惊讶，也不要认为我们所说的难以置信，因为不久以前，甚至可以说不止是以前，罗马元老院还将那些自始以来统治过他们的皇帝，或是全部，或是他们所选择与决定的，列入众神之列，并颁令崇拜他们。因为对于那些他们所敌视的皇帝，他们视之为仇敌并称之为凡人，承认其真实本性；而对于那些受他们欢迎的，则因其功绩而下令崇拜，仿佛他们自己有能力造神，尽管他们自己也是人，且并不自认超乎凡人之列。 6. 然而，如果他们能做神，他们自己理应是神；因为制造者必然优于所造之物，审判者必然高于受审判者，而施予者，至少是将自己所有的施予，乃是一种恩赐；正如每一位国王，在施予其所施的恩惠时，地位高于接受者。那么，如果他们随心所欲地定何人为神，他们自己首先应是神。但奇怪的是，他们自己作为人而死亡，恰恰揭露了他们对自己所神化之人投票的虚妄。\par

10. \Emph{希腊众神同样源出于人，由\Person{特修斯}下令。凡人成为神的过程。}\par

但这种风俗并非新近才有，也不是始于罗马元老院；相反，它自古就已存在，先前就被用于捏造偶像。因为在希腊人中久负盛名的那些神明，如\OriginalTerm{宙斯}{Zeus}、\OriginalTerm{波塞冬}{Poseidon}、\OriginalTerm{阿波罗}{Apollo}、\OriginalTerm{赫菲斯托斯}{Hephæstus}、\OriginalTerm{赫耳墨斯}{Hermes}，以及女神中的\OriginalTerm{赫拉}{Hera}、\OriginalTerm{得墨忒耳}{Demeter}、\OriginalTerm{雅典娜}{Athena}和\OriginalTerm{阿耳忒弥斯}{Artemis}，都是根据希腊历史所述，由\OriginalTerm{忒修斯}{Theseus}下令被授予神号的；并且那些颁布此类命令的人像凡人一样死去并受人哀悼，而那些被授予神号的人却被当作神明来敬拜。这是何等的不一致和疯狂！他们明知颁令者是谁，却对受令者更加尊敬。2. 但愿他们偶像崇拜的狂热仅限于男性，而没有将神号贬降至女性。因为连那些不适合参与公共事务商议的妇女，他们竟也如同事奉天主一般来敬拜事奉，比如上文所述特修斯所指定的那些女神，还有埃及人的\OriginalTerm{伊西斯}{Isis}、\OriginalTerm{少女神}{the Maid}和\OriginalTerm{幼女神}{the Younger one}，以及其他民族的\OriginalTerm{阿佛洛狄忒}{Aphrodite}。至于其他神明的名字，我甚至觉得不宜提及，因为它们充满了各种荒诞不经。3. 因为不仅在古代，就是在我们当代，许多人丧失了自己的亲人、兄弟、亲戚和妻子；还有许多失去丈夫的妇女——这些人的本性都证明他们是必死的人——他们为之造像、设计祭仪，并将之祝圣；而后世之人，被艺术家的形象和光彩所打动，就把他们当作神明来敬拜，从而违背了自然之理。因为他们的父母曾为他们哀哭，并不把他们当作神明（因为如果他们知道他们是神，就不会像对死者那样哀悼他们；这正是他们为之造像的原因，因为他们不仅不认为他们是神，甚至根本不相信他们继续存在，他们造像的目的只是为了借着看见形象来安慰自己，因为人已不在），然而愚昧的人却向他们祈祷，把他们当作神明，赋予他们只应归于真天主的尊荣。4. 例如，在埃及，直到今天，人们还在为\OriginalTerm{俄西里斯}{Osiris}、\OriginalTerm{荷鲁斯}{Horus}和\OriginalTerm{提丰}{Typho}及其他神祇举行哀歌。\OriginalTerm{多多纳}{Dodona}的铜锅和\OriginalTerm{克里特}{Crete}的\OriginalTerm{科律班忒斯}{Corybantes}，都证明\OriginalTerm{宙斯}{Zeus}不是神而是一个人，一个食人父亲所生的人。说来奇怪，连备受希腊人景仰的哲人\OriginalTerm{柏拉图}{Plato}，尽管他对天主有着自负的理解，却随\OriginalTerm{苏格拉底}{Socrates}下到\OriginalTerm{比雷埃夫斯}{Peiræus}去敬拜\OriginalTerm{阿耳忒弥斯}{Artemis}——一件人手所造的幻象。\par

\Emph{异教神祇的作为，尤其是\OriginalTerm{宙斯}{Zeus}的作为。}\par

对于这些以及类似的偶像崇拜狂热的捏造，\WorkTitle{圣经}早已预先教导我们说：\BibleQuote{实在，偶像是邪淫的开始，一发明出来，便成了生命的腐败。起初原无偶像，也不能长久存在；由于人的虚荣心，才误入世界；为此，必注定不久就要灭亡。有一父亲，因儿子夭折，忧伤逾恒，便为那早丧的儿子造像，后来，竟视之若神明，传令家属敬拜祭献。如此，相沿成习，日久而变成法律，君王的命令，亦造成人们崇拜偶像。人若在远处不能亲近，便设想对方的面貌，为所要尊敬的人，作一显明的肖像，以便公开炫耀那不在眼前的人，如同就在跟前一般。艺术家的野心，更把不认识的人予以神化，为讨好君王，尽力使其所代表的肖像美丽：民众看见这般精工，便把那原来受尊敬的人，竟视为神明。这就是世人的错觉，因为遭遇祸患，或受压迫的人们，便以\InnerQuote{不可道出的名号}，加于木石。}2. 偶像之发端的起源与设计既是如此——如\WorkTitle{圣经}所证实的——现在该是用证据驳斥它的时候了，这些证据更多是出自这些人自己对偶像的见解，而非外来之言。因为从最低处说起，若察考他们称为神者之所为，就会发现他们不仅不是神，甚至还是人类中最可鄙的。因为看看诗人们笔下\OriginalTerm{宙斯}{Zeus}的爱情与放纵行为该是怎样的景象！听听这些事吧：他一方面掳走\OriginalTerm{伽倪墨得斯}{Ganymede}，暗中行奸淫，另一方面又惊慌恐惧，生怕特洛伊人的城墙在他意料之外被摧毁！看看他因儿子\OriginalTerm{萨尔佩冬}{Sarpedon}之死而哀伤，想要救助他却无能为力，以及当他被其他所谓的神——即\OriginalTerm{雅典娜}{Athena}、\OriginalTerm{赫拉}{Hera}和\OriginalTerm{波塞冬}{Poseidon}——密谋反对时，竟靠一个女人\OriginalTerm{忒提斯}{Thetis}和百手\OriginalTerm{埃该翁}{Ægaeon}而得救，还有他屈服于享乐，成为女人的奴仆，并为她们的缘故投身冒险，化作各样的畜生、爬虫和飞鸟；还有，他因父亲设计谋害他而躲藏，或他捆绑了\OriginalTerm{克洛诺斯}{Cronos}，或他再次伤残自己的父亲！唉，将行此等事者视为神，且此人被控告之事，就连罗马的公法也不允许凡人去行，这难道是合适的吗？\par

\Emph{归诸异教神祇的其他可耻行为。这一切都证明他们不过是昔时之人，且非善类。}\par

为了免于冗长，仅从众多事例中略举数端：谁若看到他对待\Person{塞墨勒}、\Person{勒达}、\Person{阿尔克墨涅}、\Person{阿尔忒弥斯}、\Person{勒托}、\Person{迈亚}、\Person{欧罗巴}、\Person{达那厄}和\Person{安提俄珀}的无法无天与贪腐败坏的行为，或者看到他对自己的姐妹所胆敢做的事，竟使同一女人兼具妻子与姐妹的身份，谁不鄙夷他，判定他死有余辜呢？因为他不但犯奸淫，还将奸淫所生的子女神化并举到天上，以\OriginalTerm{神化}{deification}来遮掩他的不法：诸如\Person{狄俄尼索斯}、\Person{赫拉克勒斯}、\Person{狄俄斯库里}、\Person{赫尔墨斯}、\Person{珀尔修斯}和\Person{索忒拉}。2. 谁若看到那些所谓的神祇在\Place{特洛伊}因希腊人和特洛伊人而彼此势不两立地争斗，会不认识到他们的软弱呢？他们因彼此嫉妒，竟怂恿凡人也卷入纷争。谁若看到\Person{阿瑞斯}与\Person{阿佛洛狄忒}被\Person{狄俄墨得斯}所伤，或者\Person{赫拉}与来自地下的\Person{艾多纽斯}（他们称之为神）被\Person{赫拉克勒斯}所伤，\Person{狄俄尼索斯}被\Person{珀尔修斯}所伤，\Person{雅典娜}被\Person{阿耳卡斯}所伤，以及\Person{赫淮斯托斯}被摔下而跛行，会不认清他们的真实本相，并拒绝称他们为神，反而（当听说他们可朽坏且\OriginalTerm{能受苦的}{passible}时）确信他们不过是人，而且是软弱的人，且羡慕那施伤者而非受伤者呢？3. 谁若看到\Person{阿瑞斯}与\Person{阿佛洛狄忒}通奸，\Person{赫淮斯托斯}设计捉拿二人，而其他所谓的神祇被\Person{赫淮斯托斯}叫来观看通奸场面，且前来见了他们的放荡，会不耻笑而识破他们卑劣的品性呢？谁若见到\Person{赫拉克勒斯}对\Person{翁法勒}的醉狂胡为，会不耻笑呢？至于他们的享乐行为，无度的情欲，以及用金、银、铜、铁、石、木制成的神像，我们无需费力辩论揭露，这些事实本身即已可憎，单凭它们就足以证明骗局；因此，人主要的情感应是对那些受骗者寄予怜悯。4. 因为，他们憎恨那勾引自己妻子的奸夫，却不知羞耻地将奸淫的导师神化；自己禁绝乱伦，却崇拜那些行乱伦者；承认败坏儿童是恶事，却侍奉那些被控犯此罪者，毫不脸红地将法律禁止世人之间存留的事归于他们所谓的神。\par

\Emph{偶像崇拜的愚昧及其对艺术的玷辱。}\par

再者，当他们崇拜木石之物时，并不看见自己虽践踏焚烧那毫无二致的材料，却称这些材料的某部分为神。他们不久前还使用的东西，竟在愚昧中雕刻敬拜，毫不察觉也不思量，自己所拜的不是神，乃是雕刻匠的手艺。2. 因为，石头未经雕凿、木头未经加工时，他们践踏其一，频繁使用另一供自己之用，甚至作不体面的用途。但当艺术家赋予它们自己技艺的比例，将男人或女人的形状印于材料上，随即感谢艺术家，把它们当作神来崇拜，花钱从雕刻匠买来。而且，造像者往往仿佛忘了自己亲手所做的工，向自己的制品祈祷，称刚才还在削凿的东西为神。3. 然而，若真需要欣赏这些东西，不如将之归于巧匠的技艺，而不应重制品轻制作者。因为不是材料装饰了技艺，而是技艺装饰并神化了材料。那么，他们崇拜艺术家远比崇拜制品公道得多，一因他的存在先于艺术所造之神，二因这些神是按他喜悦的样式形成的。但如今，他们搁置正义，轻视技艺工艺，却崇拜技艺工艺的制品；制作的人死后，他们尊崇其作品为不朽，而作品若不天天加以看顾，迟早必自然终结。4. 再者，谁不因这点也怜悯他们：他们有眼能看，却崇拜看不见的；有耳能听，却向听不见的祈祷；生而具有生命与理性的人类，竟称那些完全不能动、甚至没有生命的东西为神；最奇怪的是，他们竟把自己掌管下的受造物当作主子来事奉。不要以为这只是我个人的说法，或是我在毁谤他们；这一切的验证就在眼前，凡愿意的都能见到类似情形。\par

\Emph{圣经谴责偶像崇拜。}\par

然而，关于这一切，有更好的证据是由圣经提供的；圣经预先告诉我们说：\BibleQuote{他们的偶像无非金银，不过是人手中的制造品：偶像有口，而不能言；偶像有眼，而不能看，有耳，而不能听；有鼻，而不能闻；有手，而不能动；有脚，而不能行；有喉，而不出声。铸造偶像的人，将与偶像同亡。}他们也没有逃脱先知的谴责；因为圣神在那里也驳斥他们说：\BibleQuote{他们必蒙羞，因为他们曾铸造神像，雕刻虚无之物；所有制造它们的人，都必枯干；让聋者们聚在一起，一同站立，他们都必惊惶，一同蒙羞。因为铁匠磨利铁器，用锛子加工，用鑽子制作，又用他有力的膀臂把它竖立；他必饥饿，筋疲力尽，不饮水，就疲乏困倦。木匠选好木料，拉墨线，用胶黏合，做成人的模样，人美丽的样式，安放在家中；这木料原是他从树丛中砍下的，是主所栽种的，雨水使它生长，为给人作燃料，人可以取来取暖，也可生火，烤饼；但他们把剩下的做成神像，加以叩拜；其实那木料的一半，他们早已在火中焚烧了。他们就在那焚烧的一半上烤肉吃，吃饱了，暖和了，说：\InnerQuote{我喜欢，因为我暖和了，看见了火。}剩下的一半，他却叩拜说：\InnerQuote{拯救我，因为你是我的神。}他们不知，也不明白，因为他们的眼昏暗而不能见，他们的心也不能领悟；他们心里不想，也不明白，自己曾把木料的一半焚烧了，在炭火上烤了饼，烤了肉吃，竟将剩下部分做成了可憎之物，还去叩拜它？要知道他们的心是灰尘，他们受了迷惑，没有人能救自己的灵魂。看哪！难道你们不说：\InnerQuote{我右手中岂不是有虚谎吗？}}\ 2. 那么，他们怎能不被众人断定为无神呢？连圣经也指责他们不虔敬。他们又怎能不是可怜的呢？他们如此公然被定罪，竟敬拜死物而非真理。他们有什么希望呢？他们又能有什么推脱之辞呢？他们信赖无感觉、无运动之物，竟以此取代真天主。\par

15. \Emph{诗人和艺术家在描绘中所传达的神明详情，显示出他们无生命，既非神，甚至也非正派的男女。}\par

要是艺术家们将神塑造得连形状都没有倒好，这样它们就不会如此明显地暴露其缺乏感知。因为若非这些偶像具备感官的象征——例如眼、鼻、耳、手、口，却没有任何真正感知和把握外界对象的动作，他们或许还能哄骗单纯者的认知，以为偶像有感觉。然而事实上，它们有这些又与没有一样，站着却不动，坐着却不安。因为它们并没有这些官能的真实活动，只是依照制作者的意愿，固定不动，毫无神明的标记，分明只是人手所安放的无生命之物。2. 又或者，这些假神的宣告者和先知——我指的是诗人和作家——倘若仅仅写下它们是神，而不去叙述它们的行为，从而暴露其不虔敬和可耻的生活，那倒好了。因为单凭神性之名，他们本可以窃取真理，或更确切地说，误导众人背离真理。然而实际上，通过叙述\Person{宙斯}的爱恋与不道德，其他神明对青年人的败坏，女神们的色欲与嫉妒，以及恐惧、懦弱和其他邪恶行为，他们只不过证明了，自己所叙述的不仅无关乎神，甚且无关乎正派人，相反，是在讲述放荡可耻之人的故事，与尊贵相去甚远。\par

16. \Emph{异教徒对此的辩解：及(1)\Quote{诗人该为这些无益的故事负责}。} 但是神的名字与存在，难道有更确凿的凭据吗？二者休戚与共。要么必须为这些行为辩护，要么就得放弃神的神性。而英雄们并不被认为会作出与其本性不符的行为，然而，按照这种辩解，神却被如此认为。\par

或许，对于这一切，不虔敬者会诉诸诗人特有的风格，说诗人特有的本领就是虚构不存在的事物，为取悦听众而讲述虚构的故事；因此他们编造了关于诸神的故事。然而，他们的这一托辞，比起其他的，从他们自己对这些事的看法和声称来看，更显得肤浅。2. 因为如果诗中所言是虚构和虚假的，那么连\Person{宙斯}、\Person{克洛诺斯}、\Person{赫拉}、\Person{阿瑞斯}及其他诸神的名字也必然是虚假的。或许如他们所说，连名字都是虚构的，虽然不存在\Person{宙斯}、\Person{克洛诺斯}或\Person{阿瑞斯}这样的存在，诗人们却虚构他们的存在以欺骗听众。但如果诗人们虚构了不存在之物的存在，他们又为何把这些当作真实存在来崇拜呢？3. 或者，他们又会说，名字虽非虚构，但他们给这些名字编造了虚构的行为。但这种辩解同样站不住脚。因为如果他们编造了行为，无疑也编造了与他们所赋予之行为相联的名字。或者，如果他们在名字上说真话，那么在行为上也必然说真话。尤其是，那些在故事中称这些为神的人，当然知道神应当如何行事，绝不会将人的观念归于神，就如同不会将火的属性归于水一样；因为火燃烧，而水的本性相反是冷的。4. 那么，如果行为配得上神，行这些事的必定是神；但如果那是人的行为，而且是不光彩之人的行为，诸如奸淫和上面提到的那些事，那样行事的人必定是人而非神。因为他们的行为必须与他们的本性相符，这样，从行为便可立刻认出行为者，从本性也可确定行为。因此，就如一个人讨论水和火，说明其作用时，不会说水燃烧而火冷却；或一个人谈论太阳和大地，不会说大地发光，而太阳播种草木果实；如果他这样说，便是极度疯狂；同样，那些作者，尤其是最杰出的诗人，如果他们真知道\Person{宙斯}及其他是神，就不会赋予他们这样的行为，表明他们不是神，而是人，而且是不庄重的人。5. 或者，如果他们作为诗人，说了谎话，而你在诽谤他们，那么他们为何不在英雄的勇敢上也说谎，把怯懦说成勇敢，把勇敢说成怯懦呢？因为他们按理也应该像对\Person{宙斯}和\Person{赫拉}那样，诽谤\Person{阿喀琉斯}缺乏勇气，吹嘘\Person{忒耳西忒斯}的威力，指责\Person{奥德修斯}迟钝，把\Person{涅斯托尔}说成鲁莽之人，叙述\Person{狄俄墨德斯}和\Person{赫克托耳}的懦弱行为，以及\Person{赫卡柏}的英勇事迹。因为归于诗人的虚构和虚假应当适用于所有情况。但事实上，他们对凡人保持了真实，却毫不羞愧地对他们所谓的神说谎。6. 有人可能会辩称，他们是在神的荒淫行为上说谎，但在赞颂中，当他们称\Person{宙斯}为诸神之父、至高者、奥林匹斯、在天上掌权时，并非捏造，而是说真话；这种辩解不仅我，任何人都能驳斥。因为真理将清楚显明，与他们相反，只要我们回顾先前的证据。因为他们的行为证明他们是人，而对他们的颂词却超越了人的本性。这两者彼此矛盾；因为天上的存在既不可能有此等行为，任何人也不可能设想如此行事的人是神。\par

17. \Emph{事实很可能是，那些丑闻故事是真的，而归于他们的神圣属性则是诗人们的谄媚所致。}\par

那么，除了颂词虚假且是谄媚，而所传的行为是真实的之外，我们还能得出什么结论呢？这一点，从普遍做法便可确知其真实性。因为没有人会在颂扬某人时同时指责其行为，相反，如果人们的行为可耻，他们会用颂词来赞扬，以掩盖丑闻，从而通过过度的赞美欺骗听众，隐藏他人的恶行。2. 正如一个人要颂扬某人，却因那人行为或任何个人品质上的丑闻而找不到颂扬的素材，便会以另一种方式赞美他，用不属于他的东西来奉承他；同样，那些奇妙的诗人，因所谓神的可耻行为而羞愧，便加给他们超人的称号，却不知道他们无法用超人的幻想掩盖人的行为，反而会因这些人的缺陷而表明，天主的属性与他们并不相称。3. 我倾向于认为，他们甚至不由自主地记述了诸神的情欲和行为。因为他们试图将\WorkTitle{圣经}所称的天主\OriginalTerm{不可通传的名号}{incommunicable name}和荣耀归给那些不是神而是必死之人，而他们的这一尝试既重大又不虔敬，因此，他们甚至违背己意，被真理所迫，不得不揭露这些人的情欲，使得有关他们的记录中所述的情欲，成为后世所有人的证据，证明他们并非神。\par

18. \Emph{异教徒的辩护续。（2）\Quote{众神因发明生活技艺而受崇拜。}但这是人类自然的成就，而非神圣的。而且，按照这个原则，为什么并非所有的发明者都被神化了呢？}\par

那么，持守这种迷信的人，能提供什么辩解，什么证据，来证明这些是真实的神呢？因为，根据上文所述，我们的论证已经表明他们不过是人，而且并非可敬之人。但或许他们会转向另一种论据，傲然诉诸他们所发现的、有益于生活的事物，说他们之所以视之为神，是因他们对人类有所贡献。因为据说\Person{宙斯}拥有造型术，\Person{波塞冬}拥有航海术，\Person{赫菲斯托斯}拥有锻冶术，\Person{雅典娜}拥有纺织术，\Person{阿波罗}拥有音乐术，\Person{阿尔忒弥斯}拥有狩猎术，\Person{赫拉}拥有制衣术，\Person{德墨忒尔}拥有农艺术，还有其他神祇拥有其他技艺，正如那些向我们讲述他们事迹的人所传说的那样。\newline{}2. 但人们应当将这些技艺及类似的技艺，不单归于诸神，更须归于人类共有的本性，因为人是通过观察自然而发现技艺的。因为就连日常言语也称技艺为模仿自然。如果他们在所从事的技艺上精熟，那并不能成为视他们为神的理由，反倒更应视他们为人；因为技艺并非他们所创造，而是在其中他们如同其他人一样模仿了自然。\newline{}3. 因为，按照有关人的定义，人天生具有获知的能力；若他们凭借人的理智，通过审察自身本性并加以认识，从而发现了技艺，这便毫无可怪之处。或者，他们若说发现技艺便有权被宣称为神，那么按照同样的理由，其他技艺的发现者也早就该被宣称为神了，正如先前那些被认为配得这一称号的人一样。因为\Place{腓尼基}人发明了字母，\Person{荷马}发明了史诗，\Person{埃利亚的芝诺}发明了辩证法，\Person{叙拉古的科拉克斯}发明了修辞术，\Person{阿里斯塔俄斯}发明了养蜂术，\Person{特里普托勒摩斯}发明了播种谷物，\Person{斯巴达的吕库古}和\Person{雅典的梭伦}发明了法律；而\Person{帕拉墨得斯}发现了字母的排列、数字、度量衡。此外，据我们的史家证实，其他人还传授了各种有助于人类生活的其他事物。\newline{}4. 那么，倘若技艺能造神，且正因技艺方才有雕刻的神像存在，那么按照他们的说法，那些后来发现其他技艺的人，必然也是神。或者，如果他们不认为这些人配享神的尊荣，而承认他们是人，那么前后一致的做法便是不称\Person{宙斯}、\Person{赫拉}和其他神祇为神，反倒应当相信他们也不过是人；而且尤其如此，因为他们即使在当年也算不上可敬之辈；正如他们为其塑像这件事本身，便清楚表明他们不过是人罢了。\par

19. \Emph{图像崇拜的不一致。为其辩护的论据。（1）神性必须借可见标记来表达。（2）图像是借天使向人传递超自然信息的一种手段。}\par

因为除了男人、女人，以及更为低等的、无理性的生物——各种各样的飞鸟、驯兽、野兽、爬虫，以及陆海和一切水域所出产之物——的形象之外，他们还能借雕塑赋予他们什么形状呢？因为人类陷入了情欲和快乐的荒谬之中，除了肉身的快乐和贪欲之外，再也看不见什么，既将心思停留在这些无理性的事物中，他们便想象神明的原理存在于无理性之物中，并雕刻出众多神祇，以迎合他们种种的情欲。\newline{}2. 他们那里有走兽、爬虫和飞鸟的像，正如那位神圣真宗教的解释者所说：\BibleQuote{他们在推理上成了虚妄的，他们愚昧的心便昏暗了。他们自称是明智的，反而成了愚妄的，将不朽天主的荣耀，改为可朽坏的人、飞鸟、四足走兽和爬虫形像的样式。因此，天主任凭他们陷入不洁的情欲。}因为，如我前文所说，他们先前既以快乐的荒谬感染了自己的灵魂，随后便沉沦至这种造神的地步；既然已经堕落，此后便好像因背弃天主而被抛弃了一样，他们就在这些事物中打滚，竟用无理性的形状来描绘\OriginalTerm{天主}{God}、\OriginalTerm{圣言}{Word}之\OriginalTerm{父}{Father}。\newline{}3. 对于这些，那些在希腊人中号称哲学家和博学之士的人，虽然被迫承认他们可见的神只是人及无理性之物的形象，却辩护说，他们之使用这些事物，旨在借此使神祇能回应他们，并得以彰显；因为否则，若非借这些雕像和仪式，他们将无法认识那不可见的天主。\newline{}4. 而那些声称提出比这些更为深入、更具哲学性理由的人则说，之所以准备并塑造偶像，是为了呼请和彰显神圣的天使与大能，好借这些方式显现，教导人认识天主；又说它们为人充当文字，通过参考这些，人可以从神圣天使借其显示的作为中，学习领悟天主。这便是他们的神话——因为我们绝不能称之为神学。但若仔细考察这论据，就会发现这些人的见解，与先前所述之人的见解一样，也是错误的。\par

20. \Emph{然而，图像这种被认为的效能究竟何在？在于质料，在于形式，还是在于制造者的技艺？所有这些观点都站不住脚。}\par

有人可以答复他们说，将他们带到真理的法庭前：天主如何通过这些物件作出回答或为人所知？是由于它们所构成的\OriginalTerm{质料}{matter}，还是由于它们所具有的\OriginalTerm{形式}{form}？如果是因为质料，那么在这些物件被制作之前，天主何必不通过一切质料无例外地显示自己，反而需要形式呢？他们徒然修建了庙宇，用来禁锢一块石头、一段木头或一块金子，而整个世界都充满这些物质。2. 但如果附加的形式是神性显现的原因，那么为何还需要金子和其余物质呢？天主大可以通过这些\OriginalTerm{像}{images}所模拟的实际自然动物来显示自己。因为按照同样的原则，关于天主的观念原本会更加高贵，如果祂通过活生生的动物，无论是理性的还是非理性的，来显示自己，而不是被人从无生命无运动的东西中寻求。3. 在这一点上，他们对自己犯下了极大的不虔敬。因为尽管他们厌恶并远离真实的动物、野兽、飞鸟和爬虫，或是由于它们的凶猛，或是由于它们的污秽，然而他们却将它们的形态雕刻在石头、木头或金子上，并将它们当作神明。但他们崇拜活物本身，比崇拜它们在石头中的形像要好。4. 但或许两者都不是原因，质料或形式都不是神性临在的原因，而仅仅是精巧的技艺召唤了神明，因为技艺是对自然的模仿。但如果神明因技艺而与居所内的人交流，那么既然技艺存在于人内，又何必需要质料呢？因为如果天主仅仅因技艺而显示自己，并因此这些像被当作神明来崇拜，那么崇拜并侍奉那些掌握技艺的人就是对的，因为他们也是理性的，并且自身拥有技巧。\par

21. \Emph{藉天使传达的观念更牵涉到更荒谬的不一致，即使是真的，也不证明崇拜神像是合理的。}\par

但是，对于他们所谓第二种更深刻的辩护，有人可以合理地补充如下。如果这些物件是由你们\OriginalTerm{希腊人}{Greeks}制作的，不是为了天主自己的自我显示，而是为了天使的临在，那么你们为什么将藉以呼求\OriginalTerm{威能者}{powers}的那些像，置于被呼求的威能者之上呢？因为正如你们所说，你们雕刻这些形象是为了领悟天主，却把实际的神像授予天主的尊荣和名号，从而使自己处于亵渎的地位。2. 因为你们既承认天主的威能超越这些像的渺小，并因此不敢通过它们呼求天主，而只敢呼求较低级的威能者，却自己越过这些后者，把那位你们惧怕其临在者的名号授予木石，并称它们为神，而不称其为石头和人的手工，并崇拜它们。因为即使假定它们如你们所虚称的，是供你们默观天主的\OriginalTerm{字母}{letters}，也不该给予记号比它们所表示的更大的尊荣。因为即使有人写下皇帝的名字，如果给予这文字比皇帝更大的尊荣，也不是没有危险的；相反，这样的人要受死刑；而文字是由书写者的技巧制作的。3. 同样，你们自己如果理智健全，就不会把如此大的天主性的启示贬低为物质：但你们也不会给予神像比雕刻它的人更大的尊荣。因为如果如他们所申辩的，神像如同字母一样指示天主的显现，并因此作为天主的标记值得被\OriginalTerm{神化}{deified}，那么更有理由神化那位雕刻和塑造它们的工匠，因为他比它们更强大、更神性，因为它们是按他的意愿被雕琢和塑造的。因此，如果字母值得赞美，那么书写者由于他的技艺和心智的技巧，就更远远超过它们令人惊叹。如果因此不应当把它们视作神明，那么我们就必须再次询问他们关于\OriginalTerm{偶像}{idols}崇拜的疯狂，要求他们为这些偶像具有如此形式而辩护。\par

22. \Emph{神像不能代表天主的真实形式，否则天主就成了可朽坏的。}\par

因为，若他们如此塑造神像的理由是神明具有人形，那么为何又赋予它们非理性生物之形呢？或者，若神明之形是非理性生物，为何他们又以理性生物之形来表现呢？又或者，若两者兼具，而他们认为神明是两者之结合，即具有理性与非理性之形，那为何他们要分离这结合之物，将兽与人的形象分开，而不是总是以两者之形雕刻，如同神话虚构中的\OriginalTerm{斯库拉}{Scylla}、\OriginalTerm{卡律布狄斯}{Charybdis}、\OriginalTerm{人马}{Hippocentaur}，以及\Person{埃及人}的狗头\OriginalTerm{阿努比斯}{Anubis}那样呢？因为，他们若要如此表现，要么应当仅以此两种性质的方式塑造，要么，若只有单一之形，就不应错误地以另一种之形来表现。\newline{}2. 此外，如果他们的神像是男性，为何又赋予它们女性之形？或者，若是女性，为何又伪其形好像男性？又或者，若两者混合，就不应分开，而应两者结合，照所谓\OriginalTerm{雌雄同体}{hermaphrodites}之样，好让他们的迷信给观者带来的不仅是不虔和诽谤的闹剧，而且还是嘲讽的场面。\newline{}2. 普遍而言，如果他们认为神明是有形体的，以至于他们为它设计出肚子、手和脚，还有脖子，以及乳房和构成人的其他器官，那么瞧，他们的心灵堕落到何等不虔和不信神的地步，竟对神明怀着这样的观念。因为随之而来的必然是，它也能遭受所有其他肉身的灾难，能被切割、分裂，甚至完全消灭。但这些及类似之事，并非天主的属性，而是属世之物的属性。\newline{}3. 因为，天主是无形的、不朽坏的、不死的，无需任何东西来维持；而这些偶像却是可朽坏的，且是肉身之形，需要肉身的服侍，一如我们之前所言。因为我们常见，老旧的偶像被更新，那些被时间、雨水，或地上的某些动物所损毁的，又被修复。在这方面，人们必须谴责其愚昧：他们竟将自己亲手制造之物宣称为神，向他们自己所装点、以艺术保护它们免于朽坏的物件祈求救恩，向他们清楚知道需要自己照料之物乞求供给，并且不以称那些被他们关在小房间里的受造物为天地的主宰为耻。\par

\Emph{23. 各种偶像崇拜的差异证明它们是虚假的。}\par

然而，不仅从这些考量中人们可以认识到他们的不信神，还可以从他们关于偶像本身的相反意见中看到。因为，若按照他们的宣称和设想，他们是神，那么人应当向哪一位效忠，应当判断哪一位为至高，以便或能满怀信心地崇拜天主，或者如他们所言，能毫不含糊地借他们认识神呢？因为并非在所有民族中，被称为神的都是同一位；相反，几乎每个民族都想象出一位各自的神。而且，甚至一个地区、一个城镇，也会为其偶像的迷信而在内部不和。\newline{}2. 例如，\Person{腓尼基人}不认识\Person{埃及人}中被称为神的那些，\Person{埃及人}也不敬拜\Person{腓尼基人}所拥有的那些偶像。\Person{西徐亚人}拒绝\Person{波斯人}的神，而\Person{波斯人}拒绝\Person{叙利亚人}的神。\Person{皮拉斯基人}也弃绝\Place{色雷斯}的神，而\Place{色雷斯}人不知道\Place{底比斯}人的神。此外，\Person{印度人}与\Person{阿拉伯人}在偶像上不同，\Person{阿拉伯人}与\Person{埃塞俄比亚人}不同，\Person{埃塞俄比亚人}又与\Person{阿拉伯人}不同。\Person{叙利亚人}不拜\Person{基利家人}的偶像，而\Person{卡帕多细亚}民族称其他存在为神。\Person{比提尼亚人}采纳了另外的神，\Person{亚美尼亚人}又想象出别的。我又何必多举例子呢？大陆上的人拜的偶像与岛上的人不同，而岛上的人事奉的又异于大陆的。\newline{}3. 而且，普遍而言，每个城市和村庄，不认识邻人的神，都偏爱自己的，认为唯有这些才是神。关于\Place{埃及}的可憎之物，甚至无需多言，因为它们尽人皆知：各城的宗教彼此对立、互不相容，邻居们总是刻意崇拜与邻近之人相反的偶像；以致于，被某些人祈求的鳄鱼，其邻居却视之为可憎；而被他人当作神崇拜的狮子，其邻居不但不拜，反而视之为野兽而杀之；被某些人尊为圣的鱼，在另一地却被用作食物。由此便在他们当中引发了争斗、暴乱和频发的流血事件，以及各种放纵情欲的事。\newline{}4. 说来奇怪，据历史学家所言，正是那些从\Person{埃及人}那里学到众神之名的\Person{皮拉斯基人}，却不认识\Place{埃及}的神，反而敬拜其他神。而且，一般而言，所有痴迷于偶像的民族，都有不同的意见和宗教，在任何情况下都无一致性可言。这也不足为奇。\newline{}5. 因为，他们从那独一天主的默观中堕落了，便沦落于众多不同之对象；且背离了\OriginalTerm{圣父的圣言，万有的救主基督}{the Word of the Father, Christ the Saviour of all}，自然地，他们的理智便四处游荡。正如那些背弃太阳、进入暗处的人，循着许多无路之径转圈，看不见眼前的人，却以为那些不在的在眼前，看却不见；同样，那些背离天主、灵魂变得昏暗的人，他们的心神处于游荡状态，就像酒醉不能看见的人，想象出非真实之事。\par

\Emph{24. 一地的所谓神明，在另一地被用作祭品。}\par

因此，这绝非他们真正不信天主的轻微证据。因为各城邦和民族所信奉的神众多而各不相同，一个民族的神竟被另一个民族的神当作牺牲，于是所有神都被彼此毁灭了。实际上，某些人奉为神明的，却被其他人当作祭品和奠酒献给自己的所谓神明；而某些人的祭品，反过来正是其他人的神明。所以埃及人敬拜公牛，以及幼犊阿匹斯，而其他人却将这些动物献祭给宙斯。因为即使他们不把别人所祝圣的相同动物献祭，但通过宰杀其同类，他们似乎也献上了同样的东西。利比亚人把一种羊奉为神，称为阿蒙，而在其他民族中，这种动物却被宰杀作为牺牲献给许多神。2. 印度人敬拜狄奥尼索斯，用这个名字作为酒的象征，而其他人则把酒作为奠祭献给其他神。其他民族尊崇河流和泉水，尤其是埃及人特别尊崇水，称之为神。然而其他人，甚至那些敬拜水本身的埃及人，却用水来洗去别人和自己身上的污垢，并将用过的水可耻地倒掉。埃及的整套偶像崇拜几乎全都是由其他民族神的祭品构成的，以至于这些埃及人甚至被其他人所嘲笑，因为他们竟将不是神的东西神化；而这些偶像，不仅在其他民族中，甚至在他们自己中间，也都只是赎罪祭品和牺牲而已。\par

\Emph{人祭。其荒谬。其盛行。其灾难性后果。}\par

但是，到了这个地步，有些人竟然如此不敬虔和愚昧，竟至于宰杀真人，将这些他们塑造其形象的神的真人献祭给他们的假神。这些可怜的人竟看不出，他们所宰杀的牺牲正是他们所制造和敬拜的神的原型，而且他们正是将人献给这些神。因为可以说，他们是用同等的献同等，甚至是用高等的献低等；因为他们用活物献给死物，用有理性的人献给无生命之物。2. 例如，称作陶里安人的斯基泰人，向他们的所谓贞女神献祭船只失事的幸存者，以及他们抓获的希腊人，就这样对同族人大行不敬，暴露了他们的神的残暴性，因为天主眷顾从危险和大海中拯救出来的人，他们却加以杀戮，几乎是在对抗天主眷顾；因为他们的野蛮性格阻碍了天主眷顾的仁慈。还有些人，在战争中获胜归来时，便将俘虏按百人分组，从每组中挑出一人，将所选之人献祭给阿瑞斯。3. 行此可憎之事的不仅是因其野蛮天性而如此的斯基泰人，相反，这种行为正是与偶像和假神相联的邪恶所生的特殊结果。因为埃及人从前曾向赫拉献上这种牺牲；腓尼基人和克里特人则曾以儿童献祭来求得克罗诺斯的悦纳。甚至古罗马人也曾以人祭来敬拜所谓的拉提亚里朱庇特，方式各异，但毫无例外地都犯下并招致了污染：他们仅因实行这些凶杀就招致了污染，同时却让这种祭品的烟雾充满庙宇，玷污了自己的神殿。4. 于是，这便成了人类无数灾祸的方便之源。因为看到他们的假神喜爱这些事，他们立刻效仿其神，犯下类似的恶行，认为效仿他们认为的优越存在乃是自己的荣耀。因此，由于杀害成人和儿童以及各种放纵，人类变得稀少了。因为几乎每个城市都充斥着各种放纵，这是其神残暴性格的结果；在偶像的庙宇中，没有一个人过清醒的生活，除非那个其放纵行为被所有人见证的人。\par

\Emph{异教的道德腐败公认源自其神。}\par

例如，过去在腓尼基的庙宇里，妇女们曾坐在外面，将自己的肉体收入献给那里的神，以为通过淫行可以讨好他们的女神，并以此博取她的恩惠。而男子则否认自己的本性，不再愿意成为男性，竟穿上女人的服饰，以为这样可以取悦并尊崇他们所谓的神之母。但所有人生活堕落，彼此竞相作恶，正如保禄，基督的圣仆人所说\BibleRef{罗 1:26}：\BibleQuote{他们的女人把顺性之用变为逆性之用；男人也是如此，弃了与女人的顺性之用，彼此欲火中烧，男人与男人行了丑事} 2. 但是他们以这些和类似的方式行事，承认并证实了他们所谓神的生活也是如此。因为他们从宙斯学会了腐化青年和奸淫，从阿佛洛狄忒学会了淫乱，从瑞亚学会了放纵，从阿瑞斯学会了凶杀，并从其他神那里学到了其他类似的事，这些都是法律所惩罚，一切清醒的人所避开的。那么，还应该认为这些做这种事的是神吗？难道不该因他们行为放荡而视其比禽兽更无理性吗？难道适合把崇拜他们的人当作人，而不应当怜悯他们比禽兽更无理性，比无生命之物更无灵魂吗？因为如果他们思考过灵魂中的理性部分，他们就不会一头扎进这些事中，也不会否认真天主，基督的父。\par

\Emph{对民间异教的驳斥既已认为是定论，我们现在转向更高形式的自然崇拜。自然如何通过其各部分互相依存而见证天主，这种依存关系禁止我们将其中的任何一部分视为至高天主。此点将详细说明。}\par

然而，那些超越这些事并敬畏受造界的人，或许在被这些丑恶行径的揭露而羞愧后，会一同拒绝那在各方面都易受谴责和驳斥的事，但认为他们持有一种根深蒂固且无法反驳的见解，即敬拜宇宙及其各个组成部分。2. 因为他们将自夸，他们敬拜和侍奉的，不是木石、人形、无理性的禽鸟、爬虫和走兽，而是太阳、月亮和整个天空的宇宙，还有大地以及整个水界：他们会说，没有人能证明这些无论如何不具有神性，因为众人都明白，它们既不缺乏生命也不缺乏理性，甚至超越人性，因为一些居住在天上，另一些居住在地上。3. 因此，也有必要审视和考察这些观点；因为在这里，我们的论证也将发现其证据对他们同样有效。但在我们审视或开始论证之前，只凭受造界几乎向他们大声抗议，并指出天主是其创造者和工匠，就足够了；他统治受造界和万有，就是我们的主耶稣基督的圣父；那些自命的哲学家背离他，转而敬拜并神化由他而出的受造界，而这受造界本身却敬拜并宣认那位他们因它而否认的主。4. 因为如果人如此敬畏受造界的各个部分，并视之为神，他们很可能因这些部分相互依存而受到谴责；这些部分更借着对祂不可违背的服从法则，彰显并见证了圣言之父，祂也是这些部分的主和创造者，正如神圣律法所说：\BibleQuote{高天陈述天主的光荣，穹苍宣扬祂的化工。}5. 这事的证据并不隐晦，反而在所有良心清白的人眼前是足够清楚的，只要他们理解之眼并未完全失明。因为，若将受造界的各部分分别看待，单独思考每一个——比如太阳自身，月亮单独，又或大地与空气，热与寒，湿与干的本质，把它们从其相互的结合中分离出来——他必定发现，没有一个部分是自足的，全都需要彼此协助，靠相互帮助而存在。太阳随整个天空运转并被其容纳，永不能偏离自己的轨道；月亮和其他星辰则见证太阳给予它们的帮助；而大地明显没有雨水就不会产出作物，雨水若不借助云层，也不会降到地上；而云若没有空气，也绝不可能自行出现并存在。空气由上层空气加热，却是由太阳而非自己得到光照和明亮。6. 同样，井水和河流若无大地就绝不会存在；但大地并非自承，而是建立在水界之上，水界又被固定在宇宙的中心。环绕整个大地之外流动的海洋和浩瀚的大洋，被风吹动，风的力量将其冲往各处。而风也不是自生的，按那些就这主题著述者之说，它们源于空气，因上层空气比下层空气的炽热高温而生，并经由下层空气吹遍各处。7. 就构成物体本性的四种元素，即热、冷、湿、干而言，谁的理解扭曲到竟不知这些事物确实在结合中存在，但若分离并单独看待，它们根据更丰富的元素的主导能力，甚至会互相毁灭？因为热量若被更多的冷量压倒便被消灭，冷又因热力而被驱散，干被湿润，湿又被干燥。\par

28. \Emph{然而，宇宙有机体也不可能是天主。因为那会使天主由不同部分组成，并使他可能分解。}\par

那么，这些事物如何能是天主呢，既然它们彼此需要帮助？或者，当它们自己也相互求助时，向它们祈求任何事物又怎会合适呢？因为，关于天主，公认的真理是：祂一无所缺，而是\OriginalTerm{自足且自存的}{self-sufficient and self-contained}，万有在祂内存在，祂服务万有而非万有服务祂；既然如此，宣称太阳、月亮以及其他受造部分——它们并非如此，甚至还需要彼此帮助——为天主，又怎能算为合理呢？2. 但或许，若将它们分开、单独考虑，我们的对手自己也会承认它们是依赖性的，这证明是显而易见的。然而，他们会把一切集合在一起，构成一个单一的身体，并说整体就是天主。因为，一旦整体被组合起来，就不再需要外来的帮助，整体便足以自存并在所有方面都独立自主；至少，那些自诩哲学家的人会这样告诉我们，只是在这里将再次被驳倒。3. 现在这个论证，丝毫不亚于先前所论及的，将显明他们的不敬虔与极度的无知。因为，如果部分的组合构成整体，而整体由部分组成，那么整体便由各部分组成，而每一部分都是整体的一部分。但这与天主的概念相去甚远。因为天主是一个整体，而非若干部分，祂不由不同的元素组成，祂自己乃是\OriginalTerm{宇宙系统的创造者}{Maker of the system of the universe}。请看他们这样说时，对\OriginalTerm{天主性}{Deity}发表了何等不敬的言论。因为，如果祂由部分组成，就必定会推论出，祂与自己不同，且由不同的部分构成。因为，如果祂是太阳，就不是月亮；如果祂是月亮，就不是大地；如果祂是大地，就不能是海洋：如此逐一考察各部分，便可发现他们理论的荒谬。4. 然而，以下从观察我们人体所得的观点，便足以驳倒他们。正如眼睛不是听觉，后者也不是手；肚子不是胸膛，脖子也不是脚，但每一肢体各有其功能，而这不同的部分构成一个单一的身体——各部分结合以发挥功用，但注定在时间进程中，当那将它们聚合在一起的本性，在天主如此安排下，按祂的旨意将它们分开时，它们便要分离——这样（但愿\OriginalTerm{天主}{He that is above}宽恕这个论证），如果他们将近受造界的各部分组合成一个身体并宣布它为天主，那么，首先，如上所述，祂便与自己不同；其次，祂注定要再次分裂，按照各部分趋于分离的自然倾向。\par

29. \Emph{自然中各力量的平衡表明，它并非天主，无论从整体还是部分来看。}\par

此外，还可以通过真理之光以另一种方式驳斥他们的无神行为。因为，如果天主本性上本是\OriginalTerm{无形、不可见、不可触摸的}{incorporeal and invisible and intangible}，他们又如何想象天主是身体，并以神的尊敬崇拜那些我们亲眼可见、亲手可触的事物呢？2. 再者，如果关于天主所说的话是真实的，即祂是全能的，没有谁的权能在祂之上，祂对万有拥有权力和统治，那么，那些将受造界神化的人何以看不出，受造界并不符合天主的这一定义呢？因为，当太阳在地底下时，地球的阴影使它的光不见，而白昼太阳又以自身光芒遮掩了月亮。冰雹常损伤地上的果实，而火若遇大水泛滥便被熄灭。春天迫使冬天退让，夏天则不容春天久留越界，而夏天自身也被秋天禁止超出其季节。3. 如果他们是天主，他们就不应彼此战胜、掩盖，而应永远共存，并同时履行各自的职能。无论黑夜白昼，太阳、月亮和众星都该同样地一起闪耀，光照万有，使一切被它们照亮。春、夏、秋、冬也该并行不悖，同时运转。海洋应与泉水混合，共同供给人类饮料。平静与暴风应同时发生。火与水应一同为人类服务。因为，若它们如对手所说是天主，又不做有害的事，反而一切都为益处，就没有人会从它们受到伤害。4. 然而，如果由于它们彼此不相容，这些事皆不可能，那么，又怎能赋予这些彼此相斥、相互冲突、无法结合的事物以天主的名称，或以应归于天主的尊敬崇拜它们呢？本性上不和谐的事物，怎能赐给他人祈祷中的和平，并成为他们和谐的源头呢？因此，太阳、月亮，或任何其他受造部分，更不用说石、金或其他材料制成的雕像，或诗人寓言中的宙斯、阿波罗等，不太可能是真天主：我们的论证已显明这一点。相反，其中一些是受造部分，另一些没有生命，还有一些不过是必死的凡人。因此，对他们的崇拜和神化并非宗教的一部分，反而是引入不敬天主和种种不虔，并且显明与他们远离了对唯一真\OriginalTerm{天主}{God}的认识，即\OriginalTerm{圣父}{Father of Christ}。5. 既然这一点已得证明，而希腊人的偶像崇拜已被显明充满了各种不虔，且其引入并非为着人的益处，而是为着人生的毁灭——那么来吧，正如我们论证起先所应许的，在驳斥了谬误之后，让我们踏上真理的道路，仰望宇宙的领袖和创造者、圣父的圣言，好借着祂认识圣父，并使希腊人知道他们已离真理多远。\par

% structure: p0061 source=h2 target=section reason=relative-heading-rank; base=h2

\section{第二部分}

30. \Emph{人的灵魂是\OriginalTerm{理智的}{intellectual}，若忠于其本性，便能凭自身认识天主。}\par

我们所谈及的那些信条已被证明不过是生活的虚假向导；但真理之路意在达到那真实而真确的天主。然而，为认识并精确领悟这真理之路，除了我们自己之外，无需任何其他东西。正如天主本身在万有之上，通往祂的道路也并非遥远，或在我们之外，而是在我们内；我们能够从自身找到它，正如\Person{梅瑟}也曾教导的，他说：\BibleQuote{信德的话就在你心里。}（\BibleRef{申 30:14}）这正是\Person{救主}所宣告并确认的，祂说：\BibleQuote{天主的国就在你们内。} 2. 因为我们内心怀有信德和天主的国，就能迅速看见并领悟那宇宙君王、\Person{圣父}的救世\OriginalTerm{圣言}{Word}。不要让崇拜偶像的希腊人寻找借口，也不要让任何人自欺，声称没有这样的道路，从而为自己的不虔敬寻找托辞。 3. 因为我们所有人都已踏上这条道路，并且拥有它，即便并非所有人都愿意循此而行，反倒偏离它，走向歧途，是因那从外吸引他们的生活享乐。若有人问，这道路是什么？我说，那就是我们每个人的\OriginalTerm{灵魂}{soul}，以及寓居其内的\OriginalTerm{智力}{intelligence}。因为唯有藉此才能默观并领悟天主。 4. 除非那些不敬神的人在否认天主之后，还要拒绝承认拥有\OriginalTerm{灵魂}{soul}；这确实比他们所说的其他话更为可信，因为拥有\OriginalTerm{理智}{intellect}的人竟否认其造作者和创造者天主，这是不合常理的。因此，为了那些单纯的人，有必要简略地说明每个人类都有\OriginalTerm{灵魂}{soul}，而且是\OriginalTerm{理性灵魂}{rational soul}；尤其因为某些教派也否认这一点，认为人不过是身体的可见形式。一旦证明了这一点，他们就将从自身获得反驳偶像的更为清晰的证据。\par

31. \Emph{理性灵魂存在的证明。（1）人与禽兽的区别。（2）人具有客观思想的能力。思想之于感觉，如同乐师之于其乐器。梦境的现象证实了这一点。}\par

首先，灵魂的\OriginalTerm{理性本性}{rational nature}因其与非理性受造物的差异而得到强有力地证实。正因为如此，一般用法以此来称呼它们，因为人类就是理性的。 2. 其次，这并不是一个普通的证明：唯有人类思考自身之外的事物，对并非当场呈现的事进行推理，运用反思，并通过判断在多种推论中选择较佳者。因为非理性动物只看见当下之事，仅因眼前之物而冲动，即使后果对它们有害；而人并非仅仅因所见而冲动，而是透过思想判断眼睛所见之物。例如，他的冲动常受理智控制；他的推理又受后反思的支配。凡是真理之友，都能察觉到人类的\OriginalTerm{智力}{intelligence}与身体感官有所区别。 3. 因此，正因有区别，智力便充当感官的审判者；当感官把握其对象时，智力则加以辨别、回忆，并指示最佳者。因为眼睛的唯一功能是看，耳朵是听，口是尝，鼻孔是闻，手是触。但应该看什么、听什么，应该触摸、品尝和闻什么，这是超越感官的问题，属于灵魂和寓居其内的智力。手能握住剑刃，口能品尝毒药，但二者都不知道这些是有害的，除非\OriginalTerm{理智}{intellect}作出判断。 4. 这情形，借助一个比喻来看，就像一把制作精良的里拉琴放在一位技巧娴熟的乐师手中。因为里拉琴的弦各有其适当的音符，或高、或低、或中间、或尖锐，但若没有乐师，其音阶无法分辨，节拍也无法把握。因为只有当握琴者弹拨琴弦，并且按准每一弦的音调时，音阶才清楚，节拍才准确。同样，感觉排列在身体中就像里拉琴，当娴熟的智力主持它们时，灵魂也就能分辨并知道自己在做什么以及如何活动。 5. 但这一点唯有人类所独有，也是人类灵魂中的理性所在，藉此人便区别于禽兽，并且显示灵魂确实与身体可见的部分不同。例如，当身体躺卧于地时，人却时常想象并默观天上的事物。当身体静止、休息和睡眠时，人却内在活动，看到自身之外的事物，旅行到其他国家，四处走动，遇见熟人，并往往借此预见和预测当天将发生的事。然而，这一切除了归因于\OriginalTerm{理性灵魂}{rational soul}，还能归因于什么呢？人在理性灵魂中思考并领悟自身之外的事物。\par

32. （3）\Emph{身体不能产生这些现象；事实上，理性灵魂的活动体现在它对身体器官本能的统御上。}\par

我们补充一点，以完成我们的论证，为那些无耻地逃入否认理性的人着想。身体在本性上是会死的，人却思考不朽之事，且常常在美德要求之处，追求死亡，这是怎么回事？或者，既然身体只能暂存，人却想像永恒之事，以致轻看眼前之事，而渴望超越之事，这又如何可能？身体不可能自发地有关于自身的这类思想，也不可能思考自身之外的事物。因为它是会死的，且只能暂存。因此，那思考与身体相反、背离其本性之事的东西，必定在种类上有别。那么，除了一个\OriginalTerm{理性而不朽的灵魂}{rational and immortal soul}之外，这还能是什么呢？因为它引入更高事物的回响，不在身体之外，而在身体之内，就像音乐家在竖琴上所做的那样。2. 再者，眼睛本性地被造为能看，耳朵为能听，它们何以避开某些对象而选择其他对象呢？是谁使眼睛转离不看呢？或者是谁使耳朵不再听，这本是它的自然功能？或者是谁常常阻止味觉，这本是它的自然冲动，去品尝事物？或者是谁制止手不去进行触摸活动，这本是其自然活动，或者使嗅觉偏离其正常运作？是谁如此违背身体的自然本能而行动？或者身体如何脱离其自然进程，转向另一者的劝诫，并任由那另一者的示意所引导？唉，这些事情恰恰证明这一点：理性灵魂统御着身体。3. 因为身体甚至没有被构成能推动自身，而是随另一者的意志被带动，正如马并不自套马具，而是由其主人驱策。因此，为人类制定了实行善行、戒避恶行的法律，而对禽兽而言，恶则不被思想，不被分辨，因为它们处于理性与理解过程之外。所以我想，我们所说的已经证明了人内存在一个理性灵魂。\par

\Emph{33. 灵魂不朽。通过（1）其与身体有别，（2）其为运动之源，（3）其想象与思想超越身体之能力，得到证明。}\par

但是，灵魂受造为不朽，是教会教导中你必须知道的进一步要点，以显示偶像将如何被推翻。但我们从对身体的认识，以及身体与灵魂的差异中，将更直接地获得对此的认识。因为如果我们的论证已经证明它与身体有别，而身体在本性上是会死的，那么灵魂是不朽的，因为它并不像身体。2. 再者，正如我们已表明的，如果灵魂推动身体而不被其他事物所推动，那么灵魂的运动就是自发的，并且这种自发运动在身体被弃置于大地之后仍会继续。如果灵魂是被身体推动的，那么其推动者的分离将导致它的死亡。但如果灵魂也推动身体，那么它更推动自身。如果它是自推动的，那么它比身体存活得更久。3. 因为灵魂的运动与其生命是同一回事，正如当然，当身体运动时我们称其为活着的，当它停止运动时我们说死亡发生。但这可以从灵魂在身体内的行动一次而永远地更清楚地表明。因为即使当它与身体结合并相连，它并未被局限或等同于身体的微小尺寸，而是常常，当身体躺在床上，不动，处于类似死亡的睡眠中时，灵魂凭藉其自身的能力保持清醒，并超越身体的本性能力，虽留在身体内却仿佛离开身体旅行，想象并注视地上的事物，甚至常常与那些超越地上和形体存在的圣人和天使交谈，并因其理智的纯洁而自信地亲近他们；那么当它在\OriginalTerm{天主}{God}——那位将它们结合者——所定的时间与身体分离时，岂不更清楚地对不朽有认识吗？因为如果即使在与身体结合时，它过着身体之外的生活，那么在身体死亡后，它的生命更将继续，并因天主——那位藉着自己的\OriginalTerm{圣言}{Word}、我们的主耶稣基督如此造它的天主——而永不止息地活着。4. 这就是灵魂思考并铭记不朽与永恒之事的原因，即因为它自身是不朽的。并且，正如身体是会死的，它的感官也以会死的事物为其对象，同样，既然灵魂沉思并注视不朽的事物，那么它是不朽的，并永远活着。因为关于不朽的观念和思想从不离开灵魂，而是驻留在它内，仿佛是其内确保其不朽的燃料。这就是为什么灵魂具有观见天主的能力，并且是它通向天主的途径，不是从外接受，而是从自身获得对\OriginalTerm{天主圣言}{Word of God}的认识与领悟。\par

\Emph{34. 因此，灵魂只要除去罪的污点，就能直接认识\OriginalTerm{天主}{God}，它自身的理性本性反照\OriginalTerm{天主的圣言}{Word of God}，它正是按\OriginalTerm{圣言}{Word}的肖像受造的。但即使它不能穿透罪在其视线前所铺展的云翳，它也面对受造物为\OriginalTerm{天主}{God}所作的见证。}\par

我们重复之前所说的话：正如人否认天主，崇拜无灵魂之物，同样，他们认为自己没有理性的灵魂，便立即受到其愚昧的惩罚，即被算作无理性的受造物。因此，因着好像缺乏自己的灵魂而迷信地崇拜无灵魂的神明，他们值得怜悯和引导。2. 但如果他们声称自己有灵魂，并以理性原则为傲，且这无可厚非，为什么他们仿佛没有灵魂一样，竟敢违背理性，不按应有的方式思考，反而把自己置于比天主更高的地位呢？因为他们拥有不朽且无形的灵魂，却用可见可朽之物塑造天主的形象。或者，为什么他们离开了天主，却不归向祂呢？因为他们能够，如同他们使自己的理智背离天主，将虚无之物冒充为神一样，同样以灵魂的理智上升，回转归向天主。3. 但是他们能够回转，只要他们放下所穿上的一切情欲污秽，并坚持不懈地洗净，直到清除了所有影响他们灵魂的外来物质，能显示出它被造时的纯朴，这样他们就能够借着它看见圣父的\OriginalTerm{圣言}{Word}，他们原是按祂的肖像被造的。因为灵魂是按照天主的肖像和模样造的，正如圣经所表明的，它借天主的口说：\BibleQuote{让我们照我们的肖像，按我们的模样造人。}\BibleRef{创 1:26} 因此，当灵魂除掉遮盖它的一切罪恶污秽，只在其纯洁中保留那肖像的模样时，那么这肖像确实被完全照亮后，灵魂就如在一面镜子中看见圣父的肖像，即圣言，并借着祂达到对圣父的认识，救主就是圣父的肖像。4. 或者，如果灵魂自身的教导不足，因为外物蒙蔽了它的理智，阻碍它看见更高的事物，但仍然可能从可见之物获得对天主的认识，因为受造界，如同用文字书写，通过其秩序与和谐，大声宣告它自己的主和创造者。\par

% structure: p0072 source=h2 target=section reason=relative-heading-rank; base=h2

\section{第三部分}

35. \Emph{受造界是天主的启示；尤其在遍及整体的秩序与和谐中。}\par

因为天主既是美善而又爱慕人类，并关怀祂所造的众灵魂——祂本性无形无像、不可理解，祂的存在超越一切受造的存在，也正因为如此，人类可能错过认识祂的道路，因为他们是从无中受造，而祂却是非受造的——为此，天主借着祂自己的\OriginalTerm{圣言}{Word}，赋予了宇宙现有的秩序，以便虽然祂本性无形无像，人却至少能通过祂的化工认识祂。因为艺术家即使不在场，也常由其作品为人所知。2. 正如人们传说雕刻家\Person{菲狄亚斯}，他的艺术作品因其对称和各部分的比例，虽不在场，却向观看的人显露了\Person{菲狄亚斯}，同样，从宇宙的秩序，人应该能领悟天主是它的创造者和技师，即使眼睛没有看见祂。因为天主并没有以祂无形的本性为依归（不要让任何人以此为借口）而把自己完全不留痕迹地隐藏于人；而是如我上面所说，祂这样安排了受造界，以至于虽然祂本性无形无像，却仍能通过祂的化工为人所认识。3. 我说这话，不是凭我自己的权威，而是依据我从谈论天主的人那里学到的，其中有\Person{保禄}，他写信给罗马人这样说：\BibleQuote{自从创造世界以来，祂那看不见的，都凭所造的万物明明可见，并可领悟；}\BibleRef{罗 1:20} 而对\Place{吕考尼雅}人说，他清楚地说：\BibleQuote{我们也是人啊！与你们有同样的性情；我们只是给你们传扬福音，为叫你们离开这些虚无之物，归依生活的天主，是他创造了天地海洋和其中的一切。他在过去的世代，容忍了万民各行其道；但他并不是没有以善行为自己作证，他从天上给你们赐了雨和结实的季节，以食物和喜乐充满你们的心。}\BibleRef{宗 14:15} 4. 谁若看到天穹的旋转，太阳与月亮的运行，以及其他星辰的位置与运动——它们以相反且不同的方向运行，但在差异中却一致遵守固定的秩序——能不得出结论：这些不是自行安排的，而是有一位区别于它们的创造者安排了它们呢？或者，谁若看到太阳白天升起，月亮夜晚照耀，且月亮的亏缺和盈满历然无差地依着固定天数变化，有些星辰按其轨道运行，轨道多样复杂，另一些则无偏行地移动，能不领悟它们确有一位创造者引导它们吗？\par

36. \Emph{若考虑形成这秩序的对立力量，这便更令人惊讶。}\par

谁看到性质相反之物结合且和谐一致，例如火与寒相混，干与湿相交，而且它们并不互相冲突，反而构成仿佛同质的一体，能不推断出有一外在于这些事物者将它们结合么？谁看到冬让位于春，春让位于夏，夏让位于秋，这些事物在本性上相反（因为一个寒冷，一个炎热，一个滋养，一个毁坏），却都产生对人类有益的平衡结果——能不认识到有一位高于它们者平衡并引导着一切，即使他看不见祂？2. 谁看到云朵承托于空中，水的重量束缚于云内，能不意识到那束缚它们并如此安排者？谁看到地，本性上最重之物，固定于水上，且在本性流动的（水）上保持不动，能不认识到有一位创造并安排者，即天主？谁看到地按时生果，天降雨，河流奔涌，井泉喷发，动物从不同父母出生，并且这些事并非随时发生，而是在特定季节——总之，在彼此不同且相反的事物中，它们所遵循的平衡而一致的秩序——能不推断出有一种能力在安排和治理它们，使其按自己的旨意妥善安排？4. 因为，若让它们自行其是，它们便无法存立，也永远不能出现，因为它们在本性上彼此对立。因为水本性沉重，趋向下流，而云则轻，属于趋向上浮、上升之类。然而我们看见，水虽是沉重，却被托上云端。再者，地极重，而水相对为轻；然而更重者却支撑在更轻者之上，地不沉落，反保持不动。雄性雌性并不相同，然而它们结合为一，结果是从二者生出与它们类似的后代。简言之，寒与热相反，湿与干相斗，然而它们相合而不相争，反而一致，产生一体，以及万物之生。\par

37. \Emph{同一主题续}.\par

因此，彼此冲突、性质相反的事物，若非有一位更高之主统合其上，使之结合，它们自己是无法和解的；各种元素如同服从主人的奴隶一样，服从于祂。它们不再各依己性与邻相争，反而认识到那统合它们的主，彼此和谐一致，它们虽本性相背，却因引导者的旨意而友善相处。2. 因为，如果它们的混融合一并非由于更高权威，重者如何能与轻者相混结合，湿者与干者，圆者与直者，火与寒，海与地，日与月，星与天，气与云，它们的本性彼此殊异，又怎能相合？因为，它们之间必生大争：一个燃烧，一个制冷；重的下坠，轻的向相反方向上升；太阳发光，而空气散布黑暗；是的，就连星辰也会彼此纷争，因为有些在上，有些在下；黑夜会拒绝为白昼让路，而坚持停留，与之争斗对抗。3. 但如果真是这样，我们看到的便不是有序的宇宙，而是混乱；不是安排，而是无政府；不是系统，而是一切失序；不是谐调，而是失调。因为在普遍的争斗与冲突中，要么一切都被毁灭，要么只有占优的原理出现。即便后者，也必彰显整体的混乱，因为若被孤留，且失去其他的帮助，它必使整体失序，正如若只剩一只手或一只脚，便无法保持身体的完整。4. 因为，要是只有太阳出现，或只有月亮运行，或只有黑夜，或总是白昼，那会是什么样的宇宙？再说，要是有天而无星辰，或有星辰而无天，又会有什么样的和谐？又或只有海，或只有地而无水及受造物的其他部分，那又有什么益处？要是元素彼此相争，或只有一个得胜，而它又不足以构成物体，人或其他动物又怎能在地上出现？因为宇宙中没有任何事物能单由热、或冷、或湿、或干构成，一切都将没有安排，没有组合。而且，即便是那看似得胜的单一元素，若没有其他元素的扶助，也无法生存：因为现在每种元素都是这样生存的。\par

38. \Emph{自然秩序的和谐彰显天主的唯一性}.\par

既然到处都不是混乱而是秩序，是匀称而非不匀称，是整齐而非凌乱，并且处于完全和谐的秩序中，我们就必然要推断并领悟到那位将万物组合、凝聚并使之和谐的主宰。因为尽管肉眼不能看见祂，但从相反事物的秩序与和谐中，可以领悟到它们的统治者、安排者和君王。 2. 同样，如果我们看见一座城，由许多不同的人组成，有大小、贫富、老幼、男女，却井然有序，居民虽彼此不同，却和谐共处，富人不反对穷人，大人不反对小人，年轻人不反对老人，而是人人平等相待，共享和平——如果我们看见这样的情形，必然推断出有一位统治者在维持和谐，即使我们没有看见他；（因为混乱是无统治的标志，秩序则显示有治理的权威：正如我们看见身体各部分的彼此和谐，眼睛不与耳朵相争，手不与脚冲突，每部分都毫无分歧地完成自己的功能，我们就由此领悟到必有一个灵魂在体内管理这些肢体，尽管我们看不见它）；因此，在全宇宙的秩序与和谐中，我们必须领悟到天主是这一切的主宰，并且祂是一，而非多。 3. 所以，这种安排上的秩序，万物的和谐一致，表明那位统治者和主宰——圣言\OriginalTerm{圣言}{Logos}——不是多，而是一。因为如果受造界有不止一位统治者，这样的普遍秩序就无法维持，万物必因众多而陷入混乱，每个统治者都试图按自己的意愿偏转整体，彼此相争。正如我们曾说过，多神论\OriginalTerm{多神论}{Polytheism}就是无神论\OriginalTerm{无神论}{Atheism}，同理，多者的统治即等于无统治。因为每一个都会取消另一个的统治，结果没有一位统治者显现，到处都是一片无政府状态。而没有统治者之处，混乱自然随之而来。 4. 反过来说，众多且不同事物的单一秩序与和谐，也表明统治者是唯一的。正如从远处听到一把七弦琴，由许多不同的琴弦组成，人们惊叹其交响的和谐，因为其乐音并非仅由低音，或仅由高音和中音组成，而是所有音调在均衡中融合——并由此必然会领悟到，这琴并非自己弹奏，也不是被不止一人拨动，而是有一位琴师，即使没有看见他，他凭其技艺将每根弦的声音融合为悦耳的交响。同样，整个宇宙的秩序既然完全和谐，没有高处与低处、低处与高处的相争，万物构成一个秩序，那么理应认为，一切受造界的统治者和君王是一，而非多，祂以自己的光照亮并推动万物。\par

39. \Emph{多神的不可能}。\par

因为我们决不可认为受造界有不止一位统治者和创造者：正确而真实的宗教则相信她的创造者是一，而受造界自身也清楚指示这一点。因为宇宙只有一个而非多个，这事实便是其创造者唯一的决定性证据。因为如果有多神，必然也会有多个宇宙。因为，由多位神创造一个宇宙，或一个宇宙由多位创造，都不合理，因为由此会产生荒谬。 2. 首先，如果独一的宇宙由多神创造，那就意味着创造者的软弱，因为众多贡献于单一结果；这强烈证明每一位的创造技艺不完备。因为若一位足够，众多就不会补足彼此的不足。但说天主有任何不足，不仅不虔敬，更是极度的亵渎。因为即使在人间，若一个工人无法独自、无需多人帮助而完成一件作品，我们也不会称他为完美的工人。 3. 但是，如果虽然每一位都能完成整体，却都参与工作以在成果中占一份，我们就得出可笑的结论：每一位都为了名声而工作，免得被怀疑无能。然而，再次，将虚荣归于神祇是最荒唐的。 4. 再者，如果每一位都足以创造整体，为何还需要多位呢？一位便足以自足宇宙。而且，将受造物弄成单一，而创造者众多且不同，显然是亵渎而可笑的，因为学问的格言是：统一而完整之物高于繁杂之物。 5. 你们必须知道，如果宇宙由多神创造，其运动就会多样化而不一致。因为考虑到每一位创造者，其运动就会相应地不同。而这种不同，如前所述，会带来凌乱和普遍的混乱；因为即使一艘船如果由多人掌舵也无法正确航行，除非一位舵手执掌舵柄；七弦琴由多人弹奏也不会发出悦耳声音，除非有一位琴师弹奏。 6. 因此，受造界为一，宇宙为一，其秩序为一，我们必须领悟它的君王和创造者也为一。正因为此，创造者自己使整个宇宙成为一，免得因多个并存而令人假设多位创造者；但正如工程是一，其创造者也应被信为一。并且，从创造者的唯一并不必然推出宇宙必须唯一，因为天主本可也创造其他宇宙。但因被造的宇宙是一，就必须相信其创造者也为一。\par

40. \Emph{宇宙的理性与秩序证明它是天主的\OriginalTerm{理性}{Reason}或\OriginalTerm{圣言}{Logos}的工程}。\par

这位创造者可能是谁呢？因为这是最需要阐明的一点，免得由于对此无知，人会假定错误的创造者，再次陷入同样古老的不敬神的错误；但我想无人对此真正存疑。因为如果我们的论证已经证明诗人们的神并非神，并且判定那些神化受造物的人为错误，且总体表明异教徒的偶像崇拜是不敬神和亵渎的行为，那么，排除这些之后，必然得出真宗教在我们这里，我们所敬拜和宣讲的天主是唯一真天主，祂是创造之主、万有的创造者。2. 那么，这位是谁呢？岂不是基督的圣父——至圣且超越一切受造存在者，祂如同一位优秀的舵手，藉着自己的智慧与自己的圣言，即我们的主、救主基督，引导、保存和秩序万物，并按照自己所见的最佳而行吗？而最佳的就是那已成就的，及我们所见正发生的，因为那是祂所愿意的；这一点人几乎无法拒绝相信。3. 因为如果受造界的运动是无理性的，宇宙漫无计划地运行，一个人或许可以合理地不相信我们所说的。但如果它存在于理性、智慧和技艺中，且全然完美地秩序，那么，那在其上并秩序它的，无非是天主的〔理智或〕圣言。4. 然而，我所谓的圣言，并非那蕴含和内在一切受造物中的、有人惯常称之为\OriginalTerm{种性原则}{seminal principle}的东西，它无灵魂，无理性或思想之能力，仅以外在技巧，凭施用者的技术运作——也不像属于理性受造物、由音节构成、以空气为表达载体的言语——我指的是那生活而大能的圣言，属于美善的天主，宇宙的天主，即那正是天主的圣言\BibleRef{若 1:1}，祂与受造之物及一切受造界不同，却是美善圣父的唯一本有圣言，藉着自己的上智安排秩序并光照这宇宙。5. 因为作为美善圣父的美善圣言，祂产生了万物的秩序，将相反之物彼此结合，使其归于一个和谐的秩序。祂作为天主的能力和天主的智慧，使天旋转，以自己的一声命悬起大地，使之稳固，虽无所凭依。受祂光照，太阳给予世界光明，月亮有定时的光辉。因着祂，水悬于云中；雨水降于大地，海洋被保持在界限内，同时大地生出青草，并披上各种植物。6. 如果有人不信地询问我们所说的，是否真有天主的圣言，那么怀疑天主的圣言的人委实疯狂；但从所见之物可以得到证明，因为万物藉天主的圣言和智慧而存在，任何受造物若非由理性——如我们所说，那理性就是天主的圣言——所造，便不会有固定的存在。\par

41. \Emph{圣言在自然界中的存在，不仅为其最初的创造，亦为其持久存有，实属必要}。\par

然而，祂虽是圣言，却如我们所说，并非像人的言语那样由音节组成；祂是祂自己圣父的不变肖像。因为人由部分构成，从虚无中造出，其言谈是复合且可分的。但天主拥有真实存在且非复合，因此祂的圣言亦拥有真实存在而非复合，而是唯一独生的天主，出于美善泉源般的圣父，在其美善中发出，秩序万有并维系之。2. 但圣言，天主的圣言，将自己与受造物结合的理由确实奇妙，教导我们事物现有的秩序恰如其分，而非其他。因为受造物的本性，既是从虚无中造出，若仅凭自身构成，则短暂、软弱且可朽。但万有之天主本性美善且极为尊高——因而慈善。因为美善者不会嫉妒任何事物：因此祂甚至不嫉妒存在，而是渴望万有存在，作为祂慈爱的对象。3. 那么，见整个受造界，就其自身法则而言，是短暂且易朽的，为使不至如此，且宇宙不至重归虚无，为此祂藉着自己永恒的圣言创造了万物，赋予受造界实质性的存在，且未让其随自身本性的漂荡于风暴中，免得它再度面临不存在的危险；而是，因为祂美善，祂以自己的圣言引导并安定整个受造界，这位圣言自己也是天主，好使藉着圣言的治理、上智安排和秩列行动，受造界能享有光明，并能永远稳固地存立。它分受了那由父获得真实存在的圣言，并藉祂获取存在，以免那若非圣言维持便会临到它的事——即朽坏——发生，\BibleQuote{因为他是不可见的天主的肖像，一切受造物的首生者，因为万物都在他内而存在，可见的与不可见的；他又是教会的头}，\BibleRef{哥 1:15-18} 正如真理的执事们在其圣著中所教导的。\par

42. \Emph{圣言的这一功能详述如下}。\par

因此，圣父的\OriginalTerm{圣言}{Logos}，全能而全美，与宇宙结合，在处处彰显祂的大能，照亮一切可见与不可见之物，将万物聚合起来，与祂自己相连，未留下任何事物不在祂的能力之内；相反地，祂使万物处处各别与整体都获得生命并得以维系；同时，祂将一切可感存在的基本元素——即热、冷、湿、干——融合为一，使其不相冲突，反而组成一首和谐的交响曲。2. 因着祂和祂的能力，火不与冷斗争，湿也不与干冲突，这些相互对立的原则，却如同朋友和兄弟一般结合起来，赋予我们所见的万物以生命，并构成物体存在的基础。地上的万物和天上的万象都听从祂，即天主圣言，因而拥有生命和秩序。因着祂，所有的海洋和大洋都在其界限内运行，正如我们上面所说，陆地长出青草，被各式各样的植物所覆盖。而且，不必花费时间列举细节，因为真理显而易见：没有任何事物存在或发生，不是由祂所造并靠祂而立，正如那神圣的经文\BibleRef{若 1:1}所言：\BibleQuote{在起初已有圣言，圣言与天主同在，圣言就是天主。万物是藉着祂而造成的；凡受造的，没有一样不是由祂而造成的。}3. 正如一位音乐家调好竖琴，运用技艺将高音与低音，以及中音与其他音调配起来，从而奏出一首和谐的曲调；同样，\OriginalTerm{天主的智慧}{Sophia}将宇宙当作一把竖琴，把空气中的事物与地上的事物，天上的事物与空气中的事物相协调，将各部分组合成整体，并以祂的意愿和命令推动万物，从而美妙而合宜地产生出宇宙及其秩序的统一。祂自身与圣父一同保持不动，却以祂的组织行动推动万物，按照祂圣父认为各物所适宜的方式行事。4. 祂神性中令人惊奇的是，祂以同一个意愿行动，同时而非间断地推动万物，整体地推动一切：无论是直的、弯的，上方的、下方的、中间的，湿的、冷的、热的、可见的和不可见的，并按照各自的本性安排它们。因为当祂一点头，直的就按直的运动，弯的也按弯的，中间的按自己的方式运动；热的获得热量，干的获得干燥，万物依照各自的本性由祂赋予生命和组织，结果产生出一首奇妙而真正神圣的和谐乐曲。\par

\Emph{说明圣言与宇宙关系的三个比喻}\par

为使如此重大的事理通过一个比喻得以明了，我们可以将所描述的比作一个大合唱团。合唱团由不同的人组成：有儿童，也有妇女，还有老人，以及年轻人。当指挥发出信号时，每个人按其本性和能力出声：男人像男人，孩子像孩子，老人像老人，年轻人像年轻人，而所有人都合成一种和谐。2. 或者，就像我们的灵魂在同一时间内，按照各个感官特有的功能推动它们，以致当某一对象出现时，所有感官便一齐发动：眼睛看，耳朵听，手触摸，嗅觉闻气味，味觉品尝——而且身体的其他部分也往往活动，例如双脚行走。3. 或者，为用第三个例子使我们的意思更明白，这就好比一座非常大的城市，是在那位建造它的统治者与君王亲自临在之下建造和管理起来的；因为当他临在并发布命令，且注视一切时，众人无不服从：有的忙于农作，有的赶快去水渠取水，另一人出去采购食物——这人去元老院，那人进议会，法官走上审判席，长官进入法庭。工匠同样从事他的手艺，水手下到海里，木匠走进工场，医生进行治疗，建筑师忙于建筑；当一人往乡下去，另一人从乡下回来，有人城里散步，有人出城又回来：但这一切都在那位唯一统治者的临在和管理之下有序进行。4. 因此，我们也必须以此方式设想整个受造界，尽管这比喻并不完全，但仍可加以扩充来想象。因为藉着天主圣言仿佛只是点头一动，万物便同时各就其位，各司其职，共同组成一个单一的秩序。\par

\Emph{将比喻应用于整个宇宙，包括可见与不可见之事}\par

因为藉着一个示意，并凭那统御并主宰万有的圣父之\OriginalTerm{圣言}{Word}的德能，天旋转，星辰运行，太阳发光，月亮循行轨道，空气承受太阳的光，以太承受它的热，风吹起；山岭高耸，海洋波涛汹涌，海中的生物生长，大地稳定不移且结果实，人被造成、生活且又死去，而万事万物都有其生命和运动：火燃烧，水冷却，泉源涌出，河流流动，季节和时辰循环，雨降下，云充满，冰雹形成，雪和冰凝结，鸟飞行，爬物行进，水族游动，海被航行，地按时播种并生长庄稼，植物生长，有的幼嫩，有的成熟，其他的在生长中老迈衰败，有些事物消逝，其他的正产生并显明。\newline{}2. 但所有这一切，以及更多因数目众多而无法提及的，那位行奇事妙迹的工匠——天主的圣言，赋予光明和生命，凭着自己的示意推动并安排，使宇宙成为一体。祂也不将无形的权能排除于自身之外；因为祂既然也是它们的创造者，就将它们也纳入宇宙之中，以自己的示意和\OriginalTerm{天主眷顾}{providence}维系并赋予它们生命。对此不信是无法推诿的。\newline{}3. 因为正如藉着祂的天主眷顾，身体生长，理性的灵魂运动，且拥有生命和思想——这几乎不需证明，因为我们看见所发生的事——同样，那天主的同一圣言，以单一简单的示意，凭自己的大能推动并维系可见的宇宙和无形的权能，给每一份子分派其适当的功能，以致神圣的权能以更神圣的方式运行，而可见的事物如我们所见那样运行。但祂自己居于万有之上，既是统治者和君王，又是组织力量，祂为了自己圣父的光荣和认识而行万事，以致几乎藉着祂所实现的事工教导我们并说：\BibleQuote{由受造物的伟大和美丽，人可推见它们的创造者}\BibleRef{智 13:5}。\par

45. 结论。\Emph{\WorkTitle{圣经}关于第一部分主题的教导}\par

因为正如举目望天，看到它的秩序和星辰的光明，便可推知那安排这些的圣言；同样，注视天主的圣言，人必然也注视到天主祂的圣父，由祂所出，故祂堪当称为圣父的解释者和使者。\newline{}2. 这也可从我们的自身经验看出；因为如果从人发出的言语，我们推知心为其源头，并且，通过思考这言语，我们用理性看出它所启示的心，那么，注视圣言的德能，我们便以更大的证据和无与伦比地更知晓祂美善的圣父，正如救主亲自说：\BibleQuote{谁看见了我，就是看见了父}\BibleRef{若 14:9}。\newline{}3. 但这一切受默感的圣经也更清楚地、更有权威地教导，致使我们如今也大胆地写信给你们，你们若查阅它们，便能证实我们所说的。\newline{}4. 因为一个论证，若由更高权威证实，便是无可反驳地确证了。因此，圣言从起初便坚定地教导犹太人民关于废除偶像，说\BibleRef{出 20:4}：\BibleQuote{你不可为你制造雕刻的偶像，也不可仿造天上、地上、或地下水中之物的形像。}而废除偶像的原因，另一位作者说明说：\Quote{外邦人的偶像是金银，是人手所造的：有口而不能言，有眼而不能看，有耳而不能听，有鼻而不能闻，有手而不能动，有脚而不能行。}创造论也没有被忽略；相反，圣经深知其美丽，免得有人只专注于这美丽而崇拜受造物，好像它们是神，而不是天主的工程，便殷切地预先教导人说\BibleRef{申 4:19}：\BibleQuote{当你举目望天，看见日月星辰和天上的众星宿时，切莫为之勾引而去敬拜事奉，因为那是上主你的天主分赐给普天下万民享用的。}但祂赐给他们，并非作他们的神，而是如我们所说，藉着这些受造物的媒介，使外邦人认识创造这一切的天主。\newline{}4. 因为古时的犹太人民早有充足的教导，他们不仅从创造的工程，也从上主的圣经，得着对天主的认识。总括而言，为引人离开偶像的错谬与非理性的幻想，祂说\BibleRef{出 20:3}：\BibleQuote{除我之外，你不可有别的神。}并不是好像有别的神，祂才禁止他们有别的神，而是免得有人离开真天主，开始为自己制造那些本不是的神，如同那些在诗人和作家中被称为神，其实并不是的。而经文本身表明它们并非神，因为它说：\Quote{除我之外，你不可有别的神}，这只是指将来的事。凡指向将来的，在说话时并不存在。\par

46. \Emph{\WorkTitle{圣经}关于第三部分主题的教导}\par

那么，那废除了异教徒无神或偶像的天主圣训，难道保持沉默，任由人类全然缺乏对天主的认识吗？绝非如此：它反而预先开启他们的明悟，说道：\BibleQuote{以色列，请听！上主你的天主是唯一的天主；} 又说道，\BibleQuote{你当全心、全力爱慕上主你的天主；} 又说道，\BibleQuote{你要钦崇上主你的天主，惟独事奉他，并依附于他。} 2. 但\OriginalTerm{圣言}{Word}对万有及对万民的眷顾与安排之大能，亦由全部受默感的圣经所证实，下列经文足可确认我们的论据，那些论及天主的人说：\BibleQuote{你奠定了大地的基础，大地就稳固不移。时日依你的法令而持续。} 又说：\BibleQuote{弹琴歌颂我们的天主，他以云彩遮蔽高天，为大地预备雨水，使群山生出青草，并生出绿草为人服务，赐予牲畜食粮。} 3. 但他藉着谁赐予呢？岂不是藉着那造成万有的吗？因为对万有的眷顾自然属于那造了万有的他；而这又是谁呢？无非是天主的圣言，论及他时，另一首圣咏说：\BibleQuote{因上主的一句话，诸天造成；因他口中的气息，诸天万象形成。} 因为他告诉我们，万物是在他内并藉着他而造成的。 4. 因此，他也劝勉我们说：\BibleQuote{他一发言，万物造成；他一下令，万物生成；} 正如卓越的梅瑟在其创造记述的开端，藉着他的叙述证实我们所说的\BibleRef{创 1:20}，说道：天主说，\BibleQuote{让我们照我们的肖像，按我们的模样造人：} 因为在他进行创造天地及万物时，圣父对他说\BibleRef{创 1:6}，\BibleQuote{让天造成，} 及 \BibleQuote{让水聚在一起，让旱地出现，} 及 \BibleQuote{让大地生出青草} 及 \BibleQuote{各种绿物：} 因此，甚至必须指摘犹太人并未真正留心圣经。 5. 因为有人可能会问，天主是向谁说话，使用命令语气呢？若他所命令和谕令的对象是他所创造的受造物，这话语便是多余的，因为它们尚未存在，而是要造成的；但没有人对不存在者说话，也不会向尚未造成的对象发出造成之命令。因为若天主对将要造成的事物下令，他必说：\Quote{天，造成吧；地，造成吧；青草，生出吧；人，受造吧。} 但事实上他并未如此；他却如此下令：\BibleQuote{让我们造人，} 及 \BibleQuote{让青草生出。} 由此证明，天主是在对身旁的某位谈论它们：既如此，当他造成万有时，有位与他同在，他正向他说话。 6. 那么，这能是谁呢，若非他的圣言？因为天主能对谁说话，除了他的圣言？或者当他创造一切受造存在时，谁与他同在，除了他的智慧，智慧说 \BibleRef{箴 8:27}：\BibleQuote{当他奠定天地时，我已在他身旁？} 但在提及天地时，天地间的一切受造物也都包括在内了。 7. 但作为他的智慧与圣言，他与他同在，望着圣父，他塑造了宇宙，组织并赋予秩序；并且，由于他是圣父的大能，他赐予万物得以存在的能力，正如救主所说：\BibleQuote{凡我看见父所作的，我也照样作。} 他的圣门徒们教导说，万有是 \BibleQuote{藉着他，并归于他} 而造成的； 8. 并且，他是善者的善嗣、真子，他是圣父的大能、智慧与圣言，并非因分沾而有，亦非这些品质自外赋予他，如同那些在他内分沾并因他而成为智慧，在他内获得能力和理智的人那样；而是他就是智慧本身、圣言本身、圣父的亲生大能、真光、真真理、真正义、真德能，且实为他的确切肖像、光辉与相似。总而言之，他是圣父的全然完美之果，是独一的圣子，是圣父不变之肖像。\par

47. \Emph{若堕落本性需得复原，则回归\OriginalTerm{圣言}{Word}乃属必要。}\par

那么，谁能以数目阐述圣父，从而发现祂圣言的德能呢？因为如同祂是圣父的圣言和智慧一样，祂也迁就受造之物，为使人认识并领悟那生祂者，而成为祂的光辉、祂的生命、门、牧者、道路、君王、统治者、万有的救主、光、赋予生命者、天主眷顾。圣父既由自己生了如此良善、身为创造主的圣子，就没有把祂隐藏起来，不让受造物看见，反而藉着万有的组织与生命——这正是祂的工程——日复一日地向万有启示祂。2. 然而，在祂内并藉着祂，圣父也启示了自己，正如救主所说\BibleRef{若 14:10}：\BibleQuote{我在圣父内，圣父在我内：}由此可知，圣言在生祂者内，受生者永远与圣父同住。但事实既是如此，没有任何事物在祂以外，无论天、地及其中的一切全都仰赖祂，人却在愚昧中舍弃了对祂的认识和事奉，而去尊崇那些不存在的事物，取代了真实而真正的天主，把不存在的事物奉为神明，\BibleQuote{事奉受造之物，而不事奉创造主，}\BibleRef{罗 1:25}从而使自己陷入愚昧和不虔之中。3. 这恰如有人欣赏作品胜过工匠，对城中的公共建筑赞叹不已，却轻看建造它们的匠人；又如有人称赞一件乐器，却藐视制造并调准这件乐器的人。愚昧而眼目极有残疾！若非造船者建造，若非建筑师修筑，若非乐师制作，他们怎能认识那建筑物、船只或琴呢？4. 这样推理的人既疯狂，又超越一切疯狂；我认为，那些不认识天主、不敬拜祂圣言——我们的主耶稣基督、众人的救主——的人，心智正是如此；圣父透过祂安排、维系万有，并对宇宙施行天主眷顾。热爱基督的朋友啊，只要你心怀对祂的信德和虔敬，就应欢欣喜乐，满怀希望，因为不朽和天国正是信德与虔诚所结的果实，只要灵魂按祂的法律装饰自己。正如那些追随祂芳表的人，奖赏是永生，那些背道而驰、不行德行的人，在审判之日必遭大耻，并且危险无可宽免，因为他们虽然认识了真理的道路，行事却与知识相反。\par

\SourceRecord{Translated by Archibald Robertson.；Nicene and Post-Nicene Fathers, Second Series；Vol. 4.；Edited by Philip Schaff and Henry Wace.；Buffalo, NY: Christian Literature Publishing Co.,；1892.；<http://www.newadvent.org/fathers/2801.htm>.}
